% 本文件由 LaTeX 排版 Agent 根据有效 translation.json 自动生成。
% author=Hippolytus; work_id=0502; title=希波里图斯圣经注释残篇

\chapter{\WorkTitle{希波里图斯圣经注释残篇}}

% structure: p0001 source=h1 target=omitted reason=page-title-already-rendered-by-parent

% structure: p0002 source=h2 target=section reason=relative-heading-rank; base=h2

\section{论六日工程}

如今我们不得不详细陈述这些事，以驳斥他人的臆测。因为有人坚持主张乐园在天上，不属于受造的体系之内。但我们亲眼看见从那里流出的河流，这些河流至今仍向任何愿意察看的人敞开，因此人人都可由此断定：乐园不属于天上，而是确实被栽植于受造的体系之中。确实，它是在东方的一个地方，是一处特选之地。\par

% structure: p0004 source=h2 target=section reason=relative-heading-rank; base=h2

\section{论创世纪}

% structure: p0005 source=h3 target=subsection reason=relative-heading-rank; base=h2

\subsection{\BibleRef{创 1:5}}

\BibleQuote{有了晚上，有了早晨，是一天。}\par

\Person{希波里图斯}：祂并没有说\Quote{黑夜和白天}，而是说\Quote{一日}，以指称光的名称。祂没有说\Quote{第一日}；因为若祂说了\Quote{第一}日，祂也必须说\Quote{第二}日被造。但合理的说法不是\Quote{第一日}，而是\Quote{一日}，以便藉着说\Quote{一}，表明它返回其轨道，并保持为一，从而构成一周。\par

% structure: p0008 source=h3 target=subsection reason=relative-heading-rank; base=h2

\subsection{\BibleRef{创 1:6}}

\BibleQuote{天主说：\Quote{在水与水之间要有穹苍。}}\par

\Person{希波里图斯}：第一天，天主从无中造了祂所造的一切。但在其他日子，祂并非从无中创造，而是从第一天所造之物中，按祂的喜悦塑造。\par

% structure: p0011 source=h3 target=subsection reason=relative-heading-rank; base=h2

\subsection{\BibleRef{创 1:6-7}}

\BibleQuote{让它在水与水之间分开；就这样成了。天主造了穹苍；并把穹苍以下的水和穹苍以上的水分开；就这样成了。}\par

\Person{希波里图斯}：由于浩大的水势在地面上奔流，因此大地是\Quote{看不见}和\Quote{无形}的。当万有的主计划使看不见的成为可见时，祂便将三分之一的水固定在中间；又将另三分之一的水单独置于高处，藉着祂自己的权能与穹苍一同升起；其余三分之一留在下面，供人类使用和益处。\Emph{此时}有一个星号。这些词句见于希伯来文，但七十贤士译本中没有。\par

% structure: p0014 source=h3 target=subsection reason=relative-heading-rank; base=h2

\subsection{\BibleRef{创 3:8}}

\BibleQuote{他们听见了主天主在乐园中散步的声响，在晚间的凉风中。}\par

\Person{希波里图斯}：更确切地说，他们因一阵微风察觉到主的来临。因此，他们一犯罪，天主就向他们显现，使他们意识到自己的罪，并召叫他们悔改。\par

% structure: p0017 source=h3 target=subsection reason=relative-heading-rank; base=h2

\subsection{\BibleRef{创 49:3}}

\BibleQuote{勒乌本，我的长子，你是我的力量，我壮年的初果；你过于暴躁，过于顽固任性：你放纵如水，不可沸腾。}\par

\Person{阿奎拉}：勒乌本，我的长子，你是我的力量，是我忧愁的总和；你在尊荣中超越，在能力中超越：你麻木如水，不可超越。\par

\Person{西马库斯}：勒乌本，我的长子，是我痛苦的开始；过度攫取，过度燥热如水，你不再超越。\par

\Person{希波里图斯}：因为天主为祂的首生之民——从埃及出来的人民——显示了极大的力量。埃及地受到多种方式的惩罚。所谓\Quote{我的力量和壮年的初果}指的是首生之民——割损之民；正如天主曾向亚巴郎和他的后裔所作出的应许。但\Quote{过于暴躁}，因为人民对天主的服从硬着颈项。\Quote{顽固任性}，因为他们不仅硬心抗拒天主的服从，而且任性妄为，甚至攻击主。\Quote{你放纵}，因为对于我们主耶稣基督，人民向父放纵无度。但\Quote{不可沸腾}，圣神如此说，作为安慰，免得它完全沸腾而溢出——给予它救恩的希望。因为沸腾溢出之物就失落了。\par

% structure: p0022 source=h3 target=subsection reason=relative-heading-rank; base=h2

\subsection{\BibleRef{创 49:4}}

\BibleQuote{因为你上了你父亲的床。}\par

\Person{希波里图斯}：他首先提及这事件——在末日，人民将攻击父的床，即新娘——教会，意图败坏她；事实上，这事直至今日仍在发生，藉着亵渎来攻击她。\par

% structure: p0025 source=h3 target=subsection reason=relative-heading-rank; base=h2

\subsection{\BibleRef{创 49:5}}

\BibleQuote{西默盎和肋未，兄弟。}\par

\Person{希波里图斯}：因为从西默盎产生了经师，从肋未产生了司铎。因为经师和司铎自愿地成全了不义，同心合意杀害了主。\par

% structure: p0028 source=h3 target=subsection reason=relative-heading-rank; base=h2

\subsection{\BibleRef{创 49:5}}

\BibleQuote{西默盎和肋未，兄弟，自愿地成全了不义。我的灵魂不要进入他们的计谋，我的心不要参与他们的集会；因为他们在愤怒中杀了人，在激情中砍断了公牛的大腿筋。}\par

\Person{希玻律}：这是指他们将要结党反对主。而且他说的这个阴谋，对我们来说是很明显的。因为蒙福的达味歌唱说：\Quote{诸侯一同商议，反抗上主}，等等。关于这个阴谋，圣神也曾预言说：\Quote{愿我的灵魂不参与}，希望尽可能引导他们离开，免得那未来的罪行藉他们发生。\Quote{他们杀了人，又砍断了公牛}；他所说的\OriginalTerm{强壮的公牛}{strong bull}是指基督。\Quote{砍断}，因为当他被悬挂在木头上时，他们刺穿了他的筋。又说：\Quote{他们发怒砍断了公牛}。请注意用词的精细：因为\Quote{他们杀了人，又砍断了公牛}。他们杀了圣徒，他们仍死去，等待复活的时候。但如同小公牛，被砍断时便倒在地上，基督也是这样，自愿顺服于肉身的死亡，但他并未被死亡所胜。虽然作为人，他成为死者之一，但在天主性中仍活着。因为基督就是那公牛——一种强壮、洁净、奉献于神圣用途的动物。圣子是一切权能的主，没有犯罪，反而为我们奉献了自己，成为馨香的祭献献于他的天主和父。因此，让那些砍断这尊贵的公牛的人听吧：\Quote{他们的忿怒这样激烈，他们的怒气这样凶残，该受诅咒！}\BibleRef{创 49:7}但这个犹太民族竟敢夸口说砍断了公牛：\Quote{我们的手流了这血。}我想，这与那句愚妄的话没有区别：\Quote{他的血（归到我们身上）}等。\BibleRef{玛 27:25}梅瑟提及\BibleRef{申 33:8}对肋未的诅咒，或者更确切地说，因该支派后来对天主的热情，特别是丕乃哈斯的热情，将其转为祝福。但对西默盎的诅咒没有解除。于是它在行动中也实现了。因为西默盎没有像其他支派那样获得产业，而是住在犹大中间。然而他的支派得以保存，虽然人数稀少。\par

% structure: p0031 source=h3 target=subsection reason=relative-heading-rank; base=h2

\subsection{创世纪 49:11}

\Emph{将他的驴驹拴在葡萄树上，将他的母驴的驹子拴在优美的葡萄树上——葡萄树的卷须——他将在酒中洗他的衣服，在葡萄的血中洗他的袍。}\par

\Person{希玻律}：\Quote{驴驹}是指外邦人的蒙召；另一头是指受割礼者的蒙召：\Quote{一头母驴}，意味着两只驴驹属于同一信仰，即两种蒙召。一头驴驹被绑在\Quote{葡萄树}上，另一头被绑在\Quote{葡萄树卷须}上，意思是外邦人的教会与主结合，而受割礼者则与律法的旧约结合。\Quote{他将在酒中洗他的衣服}：即藉着圣神和真理之言，他要洁净肉身，这里以衣服为喻。\Quote{在葡萄的血中}，即被践踏而流血的葡萄，指的是主的肉身，他洁净整个外邦人的蒙召。\par

% structure: p0034 source=h3 target=subsection reason=relative-heading-rank; base=h2

\subsection{创世纪 49:12-15}

他的眼睛因酒而喜悦，他的牙齿白如牛乳。则步隆将住在海边，他将在船只的港口旁，他的边界直达漆冬。依撒加尔选择了美物，安居在份地的中间。他见安居为美，土地肥沃，就屈肩负劳，成了农夫。\par

\Person{希玻律}：即他的眼睛因真理之言而明亮，因为它们注视所有信从他的人。他的牙齿白如牛乳——这表示他言语的光辉力量：因此他称它们洁白，并比作牛乳，因为牛乳滋养肉身和灵魂。则步隆按解释是\Quote{\Emph{芬芳}}和\Quote{\Emph{祝福}}。\par

此后，从\Person{济利禄}处摘录一段：\par

\Person{希玻律}：再者，我认为这神秘地象征我们救主的新约圣事。\Quote{他的牙齿白如牛乳}这句话表示圣事食物的卓越和纯洁。而\Quote{他的牙齿白如牛乳}这句话，我们理解为他的话给信从他的人带来光明。\par

\Person{希玻律}：此外，说则步隆将住在海边，是预言他的领土濒临大海，且以色列将与外邦人混合，两族如同合成一群。这在福音中是明显的。\Quote{则步隆地方，纳斐塔里地方}等。你将更充分地看到他的份地之丰饶，既拥有内陆领土也拥有沿海地区。\par

\BibleQuote{他就在船只停泊的港湾旁；}那就是如同安全的海港，指基督，希望的锚。这表示外邦人的蒙召——基督的恩宠将遍及全地和海洋。因为他说，\BibleQuote{他（也）在船只停泊的港湾旁，并且要延伸到漆冬。}这段预言是关于外邦人的教会的，在福音中向我们显明：\BibleQuote{则步隆地，纳斐塔里地，沿海之路，约但河对岸，外邦人的加里肋亚；那坐在黑暗中的百姓看见了皓光。}\BibleRef{玛 4:15}因此，说则步隆将居住在海边的地区，他清楚地证实了这一点，如同说将来以色列要与外邦人混杂，两民被聚集到一个羊栈，在一位首席牧人手下，即善牧（本质上），就是基督。梅瑟祝福他时说：\BibleQuote{则步隆将喜乐。}\BibleRef{申 33:18}梅瑟预言，在分配土地时，他将享有来自陆地和海洋的丰盛美物，在天主手下。\BibleQuote{在船只停泊的港湾旁；}那就是如同安全的锚地，指基督，希望的锚。因为藉着祂的恩宠，他将从许多风暴中出来，此后被带到陆地，如同停泊在港口的船只安全无虞。此外，他说\BibleQuote{他甚至延伸到漆冬，}似乎表示，在圣神的进程中，两民之间将实现完全合一，以至于以色列血统的人将占据那些在天主眼中曾经极有罪的城市。\par

取自济利禄部分之后：\par

\Person{希玻里托}：\BibleQuote{那地是肥沃的；}那就是我们主的肉：\BibleQuote{肥美}，就是\BibleQuote{丰富}；因为流蜜与奶。地的各部分被划分给他作为产业和基业——那指的是主的道理。因为这是愉快的安息，正如祂自己说：\BibleQuote{你们所有劳苦和负重担的人，到我这里来，}\BibleRef{玛 11:28}等等。因为那些遵守诫命且不否认法律规条的人，在它们和我们的主的道理中都享受安息；这就是\BibleQuote{在份地中间}的意思。如主说：\BibleQuote{我来不是要废除法律和先知，而是要成全。}\BibleRef{玛 5:17}因为我们主在遵守诫命的事实中，并没有废除法律和先知，而是成全它们，正如祂在福音中所说的。\BibleQuote{他负起劳作的担子，成了农夫。}这是宗徒们所做的。他们从天主那里接受了能力，投身劳动，成了主的农夫，耕耘大地——即人类——用我们主的宣讲。\par

% structure: p0043 source=h3 target=subsection reason=relative-heading-rank; base=h2

\subsection{创世纪 49:16-20}

\BibleQuote{丹将审判他的子民，如同自己在以色列中也是一支派。愿丹成为一条蛇在路旁，躺在小径上，咬伤马的脚跟，骑马的人要向后坠落，等候主的救恩。加得——强盗的队伍要抢劫他，但他要从脚跟追回。阿协尔——他的食物肥美，他必供给君王的佳肴。}\par

取自济利禄、阿波利纳利和狄奥多若部分之后：——\par

\Person{希玻里托}：主被描绘为一位骑手；而\BibleQuote{脚跟}指示我们\BibleQuote{末期}。祂的\BibleQuote{坠落}表示祂的死亡；正如福音所记载：\BibleQuote{看，这孩子被立定使许多人跌倒和兴起。}\BibleRef{路 2:34}我们把\BibleQuote{强盗}理解为背叛者。除了犹太人外，没有其他背叛主的人。\BibleQuote{要抢劫他}，即要阴谋反对他。\BibleQuote{在脚跟}：那是指主对埋伏攻击祂之人的帮助。并且，\BibleQuote{在脚跟}这些词语表示主将迅速复仇。他将用脚（脚跟）武装良好，追上并抢走强盗的队伍。\par

\Person{阿奎拉}：\Quote{束上腰，他必束上自己的腰；}这意味着作为一个武装和战争的人，他将武装自己。\Quote{他必在脚跟武装；}他的意思是，加得将跟在兄弟们后面武装。因为虽然他的份地在约但河对岸，但他们（该支派的人）被命令跟在兄弟们后面武装，直到他们也得到自己的份地。或者他可能的意思是，加得的宗族将生活在强盗的伪饰中，他将纠集一伙劫掠者，这正是一个\Quote{强盗的队伍}，并跟随他们，和他们一起从事海盗行为，即抢劫。\par

然而，当律法中的影子被废除，以心神和真理的敬拜被引入之时，世界需要更大的光明。为此，蒙天主默感的门徒们最终蒙召，承受了律法师的职分。因为圣咏作者曾论及犹太人的母亲——即耶路撒冷——这样说出天主的话：\Quote{你的儿子们代替了你的父亲们}；也就是说，被称作你儿子的人获得了父亲的地位。又特别论及我们的主耶稣基督：\Quote{你将使他们统治全地}。然而，他们的权柄绝非毫无烦恼。相反，他们将经历无数的患难和困惑；宗徒的历程绝不会毫无危险，正如祂以实例所暗示的，祂说：\Quote{让（丹）存在}，意思是在丹那里将有许多迫害者，如同\Quote{一条蛇卧在路旁，在道上咬伤马的脚跟}，即给予凶猛危险的咬伤；因为蛇的咬伤通常是非常危险的。它们特别咬\Quote{脚跟}，因为\Quote{祂要踏碎你的头，你要咬伤祂的脚跟}。有些人就这样迫害圣宗徒们，甚至致使他们肉身死亡。因此我们可以说，他们的处境就像一匹马失蹄踢腿时的情形。因为在这种情况下，骑手会被摔下，摔到地上，我想，他就这样等待某个活着的人。同样，蒙天主默感的宗徒们也存留着，等待他们得救赎的时刻，那时他们将受邀进入那不可动摇的国度，基督对他们说：\Quote{你们这些蒙我父祝福的，来吧}等等。\par

再者，若有人将这些话理解为并非有埋伏攻击丹的蛇，而是丹本人埋伏攻击别人，那么我们可以说，这些话指的是经师和法利塞人，这些伪君子，他们在人民中掌有审判和教导的权力，却像蛇一样攻击基督，不敬地试图迫使祂跌倒，在祂崇高而温和的行程中，以他们的螫刺折磨祂。然而，那位骑手确实跌倒了，但祂至少是自愿跌倒，自愿承受肉身的死亡。而且，祂注定要复活，有圣父作为祂的帮助者和引导者。因为圣子是天主圣父的德能，使祂自己身体的殿重新获得生命。因此，圣父拯救了祂，因为祂作为人处于危险之中，但祂的本性是天主，祂亲自维护整个受造界，无论有形无形，使其处于良好状态。同样，蒙天主默感的保禄也论及祂说：\Quote{祂虽然由于软弱而被钉在十字架上，却由于天主的德能仍然活着。}（\BibleRef{格后 13:4}）\par

阿协尔得到了托勒买和漆冬一带的地区。因此他说：\Quote{他的粮食必富有，他必向君王献上美味。}我们视此为对我们蒙召的预表；因为\Quote{富有}意味着\Quote{丰盛}。谁的粮食是丰盛的呢？难道不是我们的吗？因为主就是我们的粮食，正如祂自己所说：\Quote{我是生命的食粮。}（\BibleRef{若 6:35}）除了我们的主耶稣基督之外，还有谁能向君王献上美味呢？——不仅向信从的外邦人，也向那些受割礼且在信仰上为首的人，即列祖、圣祖、先知，以及所有相信祂的名字和受难的人。\par

% structure: p0051 source=h3 target=subsection reason=relative-heading-rank; base=h2

\subsection{\BibleBook{}{创世纪} 49:21-26}

纳斐塔里是纤细之物，在嫩枝中显示美丽。若瑟是佳美之子；我佳美的、受嫉妒的儿子；我年幼的儿子。转向我。弓箭手们合谋攻击他，辱骂他，紧紧逼迫他。他们的弓因大能而折断，手臂的筋腱因雅各伯的大能之手而松弛。从那里，那坚固以色列者来自你的天主。我的天主帮助了你，以天上的祝福祝福了你，以富有万有的地上的祝福祝福了你，以乳房和子宫的祝福祝福了你，以你父亲的祝福和母亲的祝福祝福了你。这祝福超越了永存山岳的祝福，超越了永恒丘岭的祝福；这些祝福要临到若瑟的头上，临到他弟兄的太阳穴上，他是他们的首领。\par

依玻里托：那良善而受人嫉妒的兒子，直到今日，除了我們的主耶穌基督，還能是誰呢？祂的確是那些選擇恨祂的人所嫉妒的對象，但祂絕不會被戰勝。因為祂雖然忍受了十字架，卻作為天主復活了，踐踏了死亡，正如祂的天主聖父對祂說話，並說：\Quote{坐在我的右邊。}而且，那些盡其瘋狂與祂對抗的人也歸於無有，祂已教導了我們，當祂說：\Quote{弓箭手一同商議攻擊祂，並且辱罵祂。}因為\Quote{弓箭手}——就是百姓的首領——確實召集了他們的集會，並進行了苦毒的商議。\Quote{但他們的弓被折斷，他們手臂的筋絡因雅各伯的大能者的手而鬆弛，}也就是說，因天主父，即權能之主的手，祂也使祂的兒子在天上及地上蒙受祝福。而他（納斐塔里）被採納為關乎我們之事物的預象，正如福音所顯示的：\Quote{則步隆地與納斐塔里地，沿海的路，約但河外，} \BibleRef{瑪 4:15} 等等；以及\Quote{那坐在黑暗中的百姓，看見了光明。} \BibleRef{瑪 4:17} 而這光明除了外邦人的蒙召，即主的樹幹，也就是祂的樹木，外邦人被接枝其上而結果實，還能是什麼呢？而\Quote{使嫩枝增長美麗}這句話，表達了我們蒙召的卓越。如果\Quote{使嫩枝增長美麗}這句話，或許如人所理解的，是關乎嫩枝本身，那麼這句子仍然相當可理解。因為，藉著在德行中進步，並達到更美善的事物，\Quote{向著前面的目標直跑}，\BibleRef{斐 3:15} 根據蒙福的保祿的話語，我們不斷升向更高的美麗。我當然是指屬靈的美麗，為此，將來也可對我們說：\Quote{君王戀慕你的美麗。}\par

繼阿波利納里之後的一段話：—\par

依玻里托：預言的話語再次轉向厄瑪奴耳本身。因為，依我之見，它所意指的正是已經在上述話語中所陳述的：\Quote{使嫩枝增長美麗}。因為祂的意思是指，祂增長並長成到祂起初所是的樣子，並指示回歸到祂本性所擁有的榮耀。如果我們正確領會這事，我們應說這正是\Quote{歸還}給祂的。因為正如天主的獨生聖言，身為出自天主的天主，按經上所載，祂空了自已，自願謙卑到祂先前所不是的境地，取了這卑微的血肉，並以\Quote{僕人的形像}出現，且\Quote{順服天主聖父以至於死}，同樣，此後祂被稱為\Quote{極其高舉}；似乎因著祂的人性而幾乎不擁有它，並好像出於恩寵，\Quote{得著那超乎萬名之上的名}，\BibleRef{斐 2:7} 根據蒙福的保祿的話語。但事實上，這事並非首次\Quote{賜予}祂本性所沒有的；絕非如此。相反，我們必須理解為回歸並恢復到那起初本質且不可分地存在於祂內的。正因如此，當祂按天主旨意承擔了人性卑微的狀態時，祂說：\Quote{父啊！使我得榮耀，就是未有世界以先，我同你所有的榮耀，} \BibleRef{若 17:5} 等等。因為祂在萬世之前，在世界奠基以前，與父同存，始終擁有天主性固有的榮耀。祂也很可以理解為\Quote{最小的兒子}。因為祂在末世顯現，在聖先知們光榮可敬的隊伍之後，僅僅一次，在所有那些在祂降臨之前因卓越而被算為兒子的人之後。然而，說厄瑪奴耳是嫉妒的對象，這措詞有些令人疑惑。但祂對聖徒們來說是\Quote{嫉妒或效法的對象}，他們渴望追隨祂的腳步，使自己符合祂神聖的美麗，並以祂為行為的模範，從而贏得他們最高的榮耀。另一方面，祂在另一種意義上是\Quote{嫉妒的對象}——即對那些被宣告不愛祂的人來說，是\Quote{惡意的對象}。我指的是猶太人的首領——事實上，經師和法利塞人——他們對祂懷著苦毒的嫉妒，並以祂那不能被剝奪的榮耀為誹謗的依據，用多種方式攻擊祂。因為基督確曾使死人復活，即使他們已經發臭腐爛；祂還顯示了其他神性的徵兆。這些事本應使他們驚奇，使他們準備相信，不再疑惑。然而對他們卻非如此；他們被惡意吞噬，並在心思中忍受著苦毒的劇痛。\par

繼濟利祿之後的一段話：—\par

依玻里托：除了宗徒向我們所顯示的\Quote{第二個人，出於天的主}，還能是誰呢？在福音中，\BibleRef{瑪 21:31} 祂說那承行父旨意的是\Quote{最後的}。而\Quote{回轉到我這裡}這句話，意指祂在苦難之後升到天上的父那裡。在\Quote{他們一同商議攻擊祂，並辱罵祂}這話中，所指的除了那反對我們主的人民，還能是誰呢？至於\Quote{他們極力壓迫祂}這話，誰壓迫祂，並直到今日仍壓迫祂呢？就是那些——這些\Quote{弓箭手}——以為能與主相爭的人。但雖然他們得勝將祂處死，然而\Quote{他們的弓被大力折斷}。這清楚意味著，\Quote{在復活之後}他們的弓被大力折斷。所指的乃是百姓的首領，他們列陣攻擊祂，並彷彿磨利了他們的武器尖端。但他們未能刺中祂，雖然他們行了不法之事，甚至膽敢以野獸的方式襲擊祂。\par

\Quote{你超越了永存群山的祝福。}藉着\Quote{永存和不移的群山以及永恒的丘陵}，他指的乃是圣人，因为他们超拔于尘世，不介意那朽坏的事物，只追求天上的事物，并热切渴望升至最高的德行。因此，在基督的光荣之后，便是那些最卓越的圣祖们，他们达到了德行的最高境界。然而，他们不过是仆人；惟独主、圣子，赐予他们得以显扬的方法。因此他们也承认（这话的真实性），\BibleQuote{从祂的满盈中，我们都领受了。}\BibleRef{若 1:16}\par

\Quote{我的天主帮助了你。}这清楚表明，圣子的援助和支持并非来自其他，而唯独来自我们在天之天主和圣父。而藉着\Quote{我的天主}一词，意指圣神藉雅各伯发言。\BibleRef{创 48:3}\par

\Signature{欧瑟伯}：\Quote{手臂的筋腱。}他不能说\Quote{手}或\Quote{肩膀}；但因弓的中部宽阔处被称为\Quote{臂}，他恰当地说了\Quote{臂}。\par

\Signature{希玻里}：\Quote{乳与胎的祝福。}这意指，来自天上的真正祝福乃是圣神藉着圣言降临肉体之上的恩惠。而藉着\Quote{乳与胎}，他意指童贞玛利亚的祝福。而藉着你父亲和你母亲的祝福，他也意指圣父的祝福，即我们在教会内因我们的主耶稣基督所领受的。\par

% structure: p0062 source=h3 target=subsection reason=relative-heading-rank; base=h2

\subsection{创世纪 44:27}

\Quote{本雅明是只贪婪的狼；早晨他仍要吞噬，直到晚上他才分配食物。}\par

\Signature{希玻里}：这完全适合保禄，他属于本雅明支派。因为他年轻时是一匹贪婪的狼；但当他相信后，他\Quote{分配}食物。这也是我们的主耶稣基督的恩宠向我们显明的，即本雅明支派是最早的迫害者之一，这正是\Quote{早晨}的意义。因为撒乌耳，属于本雅明支派，迫害了达味，达味被预立为主的预像。\par

% structure: p0065 source=h2 target=section reason=relative-heading-rank; base=h2

\section{出自\WorkTitle{创世纪注释}}

% structure: p0066 source=h3 target=subsection reason=relative-heading-rank; base=h2

\subsection{创世纪 2:7}

\Quote{天主用地上的灰尘形成了人。}这有何含义？我们是否要说，按照一些人的意见，造了三个人：一个属灵的，一个属魂的，一个属土的？情况并非如此，整段叙述都是关于一个人。因为\Quote{让我们造人}这话是关于将要造的人；然后经文说\Quote{天主用地上的灰尘形成了人}，所以这叙述是关于同一个人。因为他那时说\Quote{让他被造}，而如今他\Quote{造了他}，并且叙述告诉我们，他是\Quote{如何}造他的。\par

% structure: p0068 source=h2 target=section reason=relative-heading-rank; base=h2

\section{出自圣热罗尼莫\WorkTitle{书信}第36封}

依撒格象征天主圣父；黎贝加象征圣神；厄撒乌象征第一代子民和魔鬼；雅各伯象征教会或基督。依撒格\Emph{年老}，指向世界末日；他眼睛昏花，表示信德已从世界消失，宗教之光在他面前被忽视；长子被召，表明犹太人拥有法律；父亲喜爱他的肉和野味，表示从错误中拯救人，义人试图藉着教义赢得他们（字面意思：\OriginalTerm{狩猎}{hunt for}）。天主的圣言在此是祝福的重新应许，以及对将来王国的希望，在此王国中圣人将偕同基督为王，并遵守真正的安息日。黎贝加充满圣神，因为她明白她在生产前所听闻的话：\Quote{年长的要服事年幼的。}\BibleRef{创 25:23}此外，作为圣神的象征，她偏爱雅各伯。她对她的小儿子说：\Quote{到羊群去，给我取两只小山羊来。}\BibleRef{创 27:9}这预表救主在肉身中的来临，为那些被罪罚所束缚的人成就伟大的拯救；因为确实在整部圣经中，小山羊被用作罪人的象征。他奉命带来\Quote{两只}，表示接纳两个民族：\Quote{柔嫩而美好的}意指可教的和纯洁的灵魂。厄撒乌的外袍或衣服，表示希伯来人的信德和圣经，外邦人的人民被赋予了这些。放在他手臂上的羊皮，是两族人民的罪，基督当他双手被钉在十字架上时，将这些罪同他自己一起钉在上面。依撒格问雅各伯为何来得这样快，\BibleRef{创 27:20}我们视之为赞叹信者们的快速信德。献上美味，表示令天主喜悦的祭品，即罪人的得救。进食后是祝福，他因香气而喜悦。他以清晰的声音宣告复活和国度的圆满，以及他的弟兄——在以色列中信的人如何敬拜他、服事他。因为不义对抗正义，厄撒乌被激起争斗，并诡诈地谋划杀害，心里说：\Quote{为我父亲居丧的日子近了，我要杀死我的兄弟雅各伯。}\BibleRef{创 27:41}魔鬼先前在该隐身上预表了杀害兄弟的犹太人，如今在厄撒乌身上最显明地揭露了他们，也表明谋杀的时间：\Quote{让日子，}他说，\Quote{为我父亲居丧的日子来到，我好杀死我的兄弟。}因此黎贝加——即忍耐——告诉了丈夫兄弟的阴谋；他召来雅各伯，吩咐他前往美索不达米亚，从那里娶他母亲兄弟、叙利亚人拉班的家族中一个女子为妻。因此，正如雅各伯为逃避兄弟的恶计前往美索不达米亚，同样，基督也被犹太人的不信所迫，前往加里肋亚，从那里为自己娶一个来自外邦的新娘，即祂的教会。\par

% structure: p0070 source=h2 target=section reason=relative-heading-rank; base=h2

\section{论\WorkTitle{户籍纪}，出自\Quote{巴郎的祝福}}

如今，为显明祂同时兼备天主的本性与人的本性——正如宗徒也所说：\BibleQuote{在天主与人之间的中保，就是那人基督耶稣。}\BibleRef{弟前 2:5} \BibleQuote{中保原不是为一个人的，而是为两方的。}\BibleRef{迦 3:20}——因此，基督在成为天主与人之间的中保时，必要从双方领受某种凭据，好使祂显为两个不同位格之间的中保。\par

% structure: p0072 source=h2 target=section reason=relative-heading-rank; base=h2

\section{论《列王纪》}

有人提出疑问：撒慕尔是否因那女巫的作为而复活了？倘若我们承认他复活了，我们便是宣扬虚假之事。因为，一个魔鬼怎能把灵魂唤回——我不说义人的灵魂，就是任何人的灵魂——当灵魂离去，停在无人知晓的地方时？但他说：那妇人为何惊惶，又何以异常地看见人上来？倘若她的异象并非异常，她就不会说：\BibleQuote{我看见神从地中上来。}她只召唤了一个，为何有许多上来？那么，我们该说上来的所有人的灵魂都上来了，且并非被那妇人召唤；或者所见不过是他们的幻影？然而这也不够。他进一步质问：撒乌耳如何认出（所显现的），并俯伏朝拜？实则撒乌耳并未亲眼看见，只是因妇人告知那些上来的人中有一个的形象正是他所要的，便以为是撒慕尔，就向他求问，并俯伏朝拜。魔鬼要幻化出撒慕尔的形象，对它来说并非难事，因为它是认识他的。他又说：它如何预言即将临到撒乌耳和约纳堂的灾祸？它确实预言了战争的结局，以及撒乌耳如何失败，这是它从天主对撒乌耳的义怒中推断出来的。正如一位医生，虽不精于医术，但见病人已无可救药，便能说出他的死期，尽管可能在时辰上有误；同样，魔鬼因知晓天主因撒乌耳的作为，并因他这次求问女巫的举动而发怒，便预言了他的失败和死亡，尽管对他的死期有误。\par

% structure: p0074 source=h2 target=section reason=relative-heading-rank; base=h2

\section{论《圣咏集》}

% structure: p0075 source=h3 target=subsection reason=relative-heading-rank; base=h2

\subsection{引言}

希波利托斯，罗马主教，为其《圣咏注释》所写的序言。\par

《圣咏集》在梅瑟法律之后包含了新的教义。在梅瑟著作之后，它是第二部教义之书。梅瑟和若苏厄死后，在民长之后，达味兴起，他堪当被称为救主本身之父；他首先将一种新的圣咏体裁赐给希伯来人，由此废除了梅瑟所制定的关于祭献的规条，并在天主的敬拜中引入了新歌和新式欢乐赞美；他在整个服务期间，教导了许多超越梅瑟法律的事。\par

% structure: p0078 source=h3 target=subsection reason=relative-heading-rank; base=h2

\subsection{论第二篇圣咏}

出自圣主教希波利托斯对第二篇圣咏的注释。\par

祂来到世上时，显为天主和人。在祂身上，人易于察觉：祂饥饿、疲乏、困倦、口渴、因恐惧而退避、祈祷、忧愁、睡在船上枕头上、求免苦杯、痛苦流汗、被天使加力、被犹达斯出卖、被盖法嘲笑、被黑落德藐视、被比拉多鞭打、被士兵戏弄、被犹太人钉在木头上、大声呼喊将灵魂交与圣父、低头断气、肋旁被长矛刺透、裹上殓布安葬在坟墓中、第三天由圣父复活。另一方面，祂的属神性同样显明：祂被天使朝拜、被牧童看见、被西默盎等待、被亚纳证明、被贤士探寻、被星星指示、在婚宴上以水变酒、斥责被狂风掀起波浪的海、行走在深渊上、使生来瞎眼的看见、使死了四天的拉匝禄复活、行了许多奇事、赦免罪过、将权能赐给祂的门徒。\par

% structure: p0081 source=h3 target=subsection reason=relative-heading-rank; base=h2

\subsection{论第二十二篇（或二十三篇）圣咏}

出自圣主教殉道者希波利托斯对\BibleQuote{上主是我的牧者}的注释。\par

再者，用不朽坏之木造的约柜就是救主本身。因为这预表了祂（上主）那不朽坏、不腐败的居所，它不产生罪的败坏。罪人确实这样承认：\BibleQuote{我的创伤发臭流脓，因了我的愚昧。}但上主是无罪的，就祂的人性而言，是由不朽坏之木造的；也就是说，内在是由童贞女和圣神，外在是由天主的圣言，如同用纯金包裹的约柜。\par

% structure: p0084 source=h3 target=subsection reason=relative-heading-rank; base=h2

\subsection{论第二十三篇（或二十四篇）圣咏}

祂来到天门：天使伴随祂；天堂的门是关闭的。因为祂尚未升到天上。如今祂首先以肉身升起的姿态向天上的权能显现。因此，那些作为救主和上主信使的天使向这些权能说：\BibleQuote{众门哪，抬起你们的头来；永久的门户哪，你们要被举起；荣耀的王将要进来。}\par

% structure: p0086 source=h3 target=subsection reason=relative-heading-rank; base=h2

\subsection{论第一百零九篇（或一百一十篇）圣咏}

出自同一位作者对《伟大之歌》的注释。\par

1. 那将最初由泥土所造的人从最低的地狱中解救出来，当他被死亡的锁链捆绑而迷失时；那从高天降临，将尘世所生的人高举到高处；那成为向死者传福音的宣讲者、灵魂的救赎者和被埋葬者的复活者——祂在人的失败中成为人的帮助者，并按祂的肖像显现，作为首生的圣言，在童贞玛利亚中取了第一个亚当，并亲自承担了祂；祂虽是属神的，却在母胎中认识了属土的；祂虽是永生的，却认识了在过犯中死亡的人；祂虽是属天的，却将属地者带到高处；祂虽出身高贵，却因自己的服从选择了释放奴隶；祂使那归于尘土、成为蛇的食物的人坚不可摧如金刚石，且当祂被挂在木头上时，宣告他成为战胜者的主，并如此借着木头证明自己是战胜者。\par

2. 那些如今不承认降生成人的天主圣子的人，将来必得承认祂是审判者，当那如今在祂无荣耀的身体中被轻视的祂，在祂的荣耀中降临时。\par

3. 当宗徒们在第三天来到坟墓时，他们没有找到耶稣的身体；正如以色列子民上山寻找梅瑟的坟墓，却没有找到一样。\par

% structure: p0091 source=h3 target=subsection reason=relative-heading-rank; base=h2

\subsection{论圣咏第七十七或七十八篇}

45. 祂将狗蝇降在他们中间，吞灭他们；又将蛙降下，毁灭他们。\par

46. 祂也将他们的果实交给霉烂，将他们的劳苦交给蝗虫。\par

47. 祂用冰雹摧毁他们的葡萄树，用霜冻摧毁他们的无花果树。\par

正如因生活不规律，体内可能形成致命的胆汁液体，医生通过技艺可使其成为病态呕吐，而医生本身并不负责产生人体内的病态液体——因为饮食过度产生了它，而医生的科学只是使之显现——同样，尽管可以说来自天主的痛苦报应临到那些自愿作恶的人身上，但按正确理性，应当认为这类邪恶既在我们自身中始，也在我们自身中因。因为对于无过犯生活的人，没有黑暗，没有虫子，没有\OriginalTerm{地狱}{Gehenna}，没有火，也没有任何其他这些恐怖的话语或事物；正如埃及的灾祸对希伯来人并不存在——那些用无形叮咬烦扰人的虱子，用痛苦刺伤身体的狗蝇，从天上带着冰雹降临的飓风，农夫劳碌被蝗虫吞吃，昏暗的天空，以及其他。确实，天主的计划是照料真葡萄树，毁灭埃及人，却宽恕那些将要\Quote{吃毒葡萄并喝毒蛇的致命毒液}的人。埃及的无花果树也被完全毁灭；然而，并非撒该爬上去以便看见我主的那一棵。埃及的果实被浪费，即肉性的事工，但并非圣神的果子，即爱、喜乐与平安。\BibleRef{迦 5:22}\par

48. 祂也将他们的牲畜交给冰雹，将他们的财产交给火。\par

\WorkTitle{《\Person{Symmachus}译本》}译为：\Quote{祂将他们的牲畜交给瘟疫，将他们的财产交给飞鸟。}因为遭遇了毁灭性的倾覆，他们成了食肉鸟类的猎物。但根据\WorkTitle{《七十子译本》}，其含义并非冰雹摧毁了他们的牲畜，而火摧毁了其余的财产，而是冰雹与火以非常方式一同降下，首先彻底摧毁了他们的葡萄树和无花果树——它们完全无法抵挡第一次冲击；然后摧毁了平原上放牧的牲畜；再后摧毁了一切草本植物和树木——伴随冰雹的火将它们烧尽；整个事件非常奇异，因为火与水一同奔跑并混合。经上说：\Quote{因为火在冰雹中奔跑。}因此，冰雹和火在冰雹中燃烧。达味也称牲畜和树木的果实为\Quote{财产}或\Quote{财富}。还应注意，尽管记载冰雹摧毁了每一棵草本和树木，但仍留下了一些，后来被蝗虫——在火雹之后来到——吃掉；关于蝗虫，它说它吃尽一切草本和树木的果实，就是冰雹所剩下的。从属灵意义上说，有些羊属于基督，有些属于埃及人。然而，那些曾经属于别人的羊可以成为祂的，正如拉班的羊成为雅各伯的；反之亦然。此外，雅各伯无论摒弃哪只羊，都把它交给厄撒乌。因此，要当心，免得你被发现是在耶稣的羊群中，却在送礼物给厄撒乌时被分开，并作为可弃绝的、不配属灵雅各伯的而交给厄撒乌。纯朴的人是基督的羊，天主按这话拯救他们：\Quote{主，祢保护人和牲畜。}那些愚昧地依附于无神道理的人是埃及人的羊，这些也被冰雹毁灭。埃及人所拥有的一切都被交给火，但亚巴郎的财产却给了依撒格。\par

49. 祂向他们发出祂的怒气——愤怒、恼怒和磨难，由恶天使而来的惩罚。\par

在愤怒、恼怒和磨难之下，他意指痛苦的惩罚；因为天主是没有情感的。通过愤怒，你将理解为较轻的惩罚；通过恼怒，较重的；通过磨难，最重的。天使也被称为恶的，并非因为他们本性如此或出于自己的意愿，而是因为他们有这职分，并被指派带来痛苦和苦难——因此被称为恶，是就承受这类事的人的心境而言；正如审判之日被称为恶日，因为它对罪人充满苦难和痛苦。以同样的效果，依撒意亚的话说：\Quote{我，主，造平安，也造邪恶；}\BibleRef{依 45:7} 借此意思，我维持平安，也容许战争。\par

% structure: p0100 source=h2 target=section reason=relative-heading-rank; base=h2

\section{论《箴言》}

% structure: p0101 source=h3 target=subsection reason=relative-heading-rank; base=h2

\subsection{第一片段}

摘自圣\Person{希波吕托斯}对《箴言》的注释。\par

所以，箴言是对整条生命之路有益的劝戒之辞；因为对于寻求通往天主之路的人来说，这些箴言是引导和记号，在他们因路途漫长而疲倦时使他们重新振作。此外，这些是\Person{撒罗满}的箴言，即是那位\Quote{和平者}，他实是救主基督。既然我们毫不冒犯地理解主的话，作为主的话，以便没有人因名字相似而误导我们，他告诉我们谁写了这些，以及他是哪一民族的人，好让说话者的信誉能使言论被接受，听众专心；因为那是那位\Person{撒罗满}的话，主曾对他说：\BibleQuote{我要赐给你一颗明智和聪慧的心，以致在你之前没有像你一样的，在你之后也不会兴起像你一样的}\BibleRef{列上 3:12}，以及关于他所记载的随后内容。现在他是明智父亲所生的明智儿子；因此加上了\Person{达味}的名字，由他生了\Person{撒罗满}。他自幼受教于圣经，并且不是凭抽签或强力，而是凭圣神的判断和天主的谕令获得他的统治。\newline{}\par

\Quote{认识智慧和训诲。}认识天主的智慧的人，也从祂领受训诲，并借此学习圣言的奥秘；认识真正的天上智慧的人，将容易理解这些奥秘的话语。因此他说：\Quote{明白言辞的难处；}因为圣神以奇异的语言所说的事，对于那些心向天主正直的人就成为可理解的。\newline{}\par

这些话他理解为论及犹太人民，以及他们在基督的血上的罪；因为他们以为祂只在世上拥有居所（公民身份）。\newline{}\par

他们不会简单地获得，而是继承。再者，恶人即使被高举，也只是被高举到更大的羞辱。因为正如一个人高举一个丑陋畸形的人，不是尊敬他，而是更羞辱他，使他的羞耻暴露给更多人；同样，天主高举恶人，是为了使他们的耻辱显明。因为\Person{法郎}曾被高举，但只是为了让世界作他的控告者。\newline{}\par

\BibleRef{箴 4:2}必须注意，他称法律为一份美好的礼物，是为了那不义地将礼物收入怀中的人。而抛弃法律的人就是违犯它；即他所说的法律，或他所遵守的法律。\newline{}\par

而所谓\Quote{高举（巩固）她}是什么意思？用圣善的思想围绕她；因为你需要大量的防御，因为有许多事危及这样的产业。但如果巩固她在于我们的能力，而提升天主知识的美德在于我们的能力，这些将是她的堡垒——例如，实践、研习，以及其他美德的整个链条；而遵守这些的人，就尊崇智慧；其报酬是，被高举与她同在，并在天堂的洞房里被她拥抱。\newline{}\par

异端者就是\Quote{恶人}，而违犯法律的人是\Quote{邪恶的人}，他吩咐我们不可进入他们的\Quote{道路}——即他们的行为。\newline{}\par

他\Quote{正直地看}，是指拥有没有激情的思想；而他有正确的判断，是指不为外表所激动。当他说\Quote{让你的眼睛正直地看}，他指的是灵魂的视觉；当他劝诫说\Quote{我儿，吃蜂蜜，好使它在你的口中甘甜}，他比喻性地使用\Quote{蜂蜜}，意指神圣的道理，这道理恢复灵魂的灵性知识。但智慧也拥抱灵魂；因为他说：\Quote{爱她，好使她拥抱你。}而灵魂借着她的拥抱与智慧合一，充满圣洁与纯洁。不仅如此，基督的芳香油膏也被灵魂的嗅觉所捕捉。\newline{}\par

美德居中；因此他也说，勇敢介于胆大妄为与怯懦之间。现在他提到\Quote{右}，并非指本性为右的事物，如美德，而是指因它们的快乐而在你看来为右的事物。因为快乐不只是感官的享受，也包括财富和奢侈。而\Quote{左}表示嫉妒、抢夺等。因为他说：\Quote{\OriginalTerm{玻瑞阿斯}{Boreas}是苦风，却仍被称为\Quote{右}}。因为，象征地说，他用玻瑞阿斯指明邪恶的魔鬼，由它在地上点燃一切邪恶的火焰。而它被称为\Quote{右}，因为天使被称为一个正确的（吉祥的）名字。他说，离开邪恶，天主必会照顾你的终局；因为祂必在你前面行，驱散你的仇敌，使你能平安前行。\newline{}\par

\BibleRef{箴 5:19}他借着提到动物（母鹿）也显示那快乐的纯洁；又借着瞪羚暗示妻子迅速的回应之情。而他虽知道许多诱发的事，却保护他们免受其害，并用爱的不可分解的纽带将他们结合，将恒心置于他们面前。至于其他，智慧，比喻地说，像一只鹿，能击退并粉碎异端者的蛇般教义。因此，他说，让她与你同在，像一只瞪羚，使一切美德保持新鲜。而虽然妻子与智慧在此方面并不相同，倒不如让她引导你；因为这样你必怀有善念。\newline{}\par

你或许会说：眼睛有什么害处呢？既然看的人不一定必要被引入歧途？祂向你指明：情欲是火，肉体如同衣服。后者是易受攻击的，前者是暴君。当有害之物不仅被接纳在内，且被紧紧抓住时，它不会出去，直到为自己开出一条出口。因为凡注视妇女的人，即使逃过了诱惑，也不可能不带一切色欲而离开。一个人若能守贞且无忧，又何必要自寻烦恼呢？看约伯怎么说：\Quote{我与眼立了约，决不注视处女。}\BibleRef{约 31:1} 他如此深知妄用之权能。保禄也为此缘故，常\Quote{痛击己身，使之隶属于己}。打个比方说，谁若允许不洁的思念盘踞在心中，就是在胸中保留一团火。谁若在行动上犯罪而毁坏自己的灵魂，就是在炭火上行走。\par

\OriginalTerm{cemphus}{cemphus}是一种野生海鸟，其性欲冲动如此过度，以致交配时眼睛似乎充血；它常常盲目地落入罗网或人手。因此，经上用此鸟比喻那些因无度淫欲而投身妓女的人；或是因这动物愚昧无知，因为那追求的对象也如同无知觉一般。据说这鸟非常喜爱泡沫，若航海时有人手捧泡沫，它便会落在他手上。它产卵也带着痛苦。\par

\BibleRef{箴 7:26} 你已看见她的祸害。不要等待情欲升起；因为她的死亡是永远的。至于其他，她用言语、真理的论证伤人，用罪恶杀死屈服于她的人。因为邪恶有多种形式，将愚昧人引向地狱。死亡的内室指其深处或宝库。那么，如何可能逃脱呢？\par

祂意指新耶路撒冷，或圣化了的肉身。七根柱子指圣神七重一体安息在她身上，正如依撒意亚证言说：\Quote{她宰杀了}她的\Quote{祭品}。\par

请注意，智者必须有益于众人；所以那只对自己有益的人不能称为智者。因为智慧若只把能力保留给拥有她的人，那谴责是大的。正如毒药并不伤害别的身体，只伤害服下它的那一个；同样，成为邪恶的人伤害的是他自己，而非他人。因为没有一个真正有德之人会被邪恶之人伤害。\par

正义的果实和生命树就是基督。惟有祂作为人完成了全部正义。祂用自己本有的生命像树一样结出了知识和德行的果实，凡吃这果子的，必得永生，并同亚当和所有义人一起享受乐园中的生命树。但恶人的灵魂会不合时宜地被逐出天主的临在，祂任他们留在折磨的火焰中。\par

他蒙恩宠，不是从人那里，而是从主那里。\par

谁寻求知道天主的旨意，就是求问智慧。谁对自己所学到的东西谨言慎行，就显明自己谨慎。一个人询问智慧，渴望学习关于智慧的事；而另一个人什么都不问，不仅自己不愿学习智慧，还阻止邻人学习；前者当然比后者更为谨慎。\par

关于水蛭。有三女儿被罪恶所溺爱——奸淫、谋杀和偶像崇拜。这三样不满足她，因为她不会满足。通过这些行为毁灭人，罪恶从不改变，只是不断增长。至于第四个，他继续说，从不说\Quote{够了}，意指普遍的情欲。他提到\Quote{第四个}，是指普遍情欲。因为身体是一个，却有许多肢体；同样，罪恶虽是一个，却包含许多不同情欲，借此布设罗网陷害人。因此，为教导我们这事，他使用阴府（\OriginalTerm{Hades}{Hades}）、女色、地狱（\OriginalTerm{Tartarus}{Tartarus}）和不得水满足的地作为例子。水与火永不说\Quote{够了}。坟墓（\OriginalTerm{Hades}{Hades}）绝不停止接纳恶人的灵魂；对于女人情欲的罪的迷恋，也不停止落在奸淫上，并成为灵魂的叛徒。至于塔尔塔罗斯（\OriginalTerm{Tartarus}{Tartarus}），位在悲惨黑暗之处，没有一丝光照；同样，凡在一切肉体情欲中作罪恶奴隶的人，如同不得水满足的地，永不能走向告解和重生之水洗，像水与火一样，永不说\Quote{够了}。\par

正如蛇不能在岩石上留下痕迹，同样魔鬼不能在基督的身体内找到罪。因为主说：\Quote{看，这世界的首领来了，他在我身上一无所获。}——正如船在海中航行，身后不留航迹；同样，教会处在世界中如同在海中，也不把她的希望留在地上，因为她的生命保留在天上；她在此世只是短暂停留，所以无法追踪她的航程。——正如教会不把希望留在世上，她对那带给我们一切美善的基督降生成人的希望，也没有在阴府留下死亡的痕迹。——是谁呢？难道不是那由圣神和童贞女所生的一位吗？——祂在世上重建完美的人，行奇迹，从若翰的洗礼开始，正如福音作者证言：那时耶稣开始大约三十岁。这正是祂青春鼎盛的年龄，祂周游各城各地，医治人的疾病和软弱。\par

\BibleQuote{嘲笑父亲，轻视母亲年老的眼睛。}亦即亵渎天主、轻视基督之母——天主的智慧的人——愿乌鸦从溪谷啄出他的眼睛，即不洁与邪恶的灵夺去他喜乐的明眸；愿雏鹰吞吃他；这样的人必被践踏在圣人脚下。\par

\BibleQuote{有三件事我无法理解，第四件事我不知道：鹰在天空飞行的踪迹，}即基督的升天；\BibleQuote{蛇在磐石上爬行的踪迹，}即魔鬼在基督的身体上找不到罪的痕迹；\BibleQuote{船在海上航行的踪迹，}即教会在这如海的今生中，藉着对基督的盼望经过十字架而前行；\BibleQuote{以及男子在青年时期的踪迹，}——即由圣神与童贞玛利亚所生者之踪迹。因为看，经上说：有一人，其名为\Quote{升起}。\par

\BibleQuote{淫妇的行径也是如此：她行完罪孽后，便洗净自己，还说没有行恶。}信仰基督的教会也是如此：当她与偶像行淫后，便弃绝它们和魔鬼，洗净自己的罪，获得赦免，然后声称自己未行恶。\par

\BibleQuote{大地因三样事震动，}即因圣父、圣子、圣神。\BibleQuote{第四样它不能承受，}即基督的最后显现。\BibleQuote{仆人作王时：}以色列在埃及为奴，在应许之地成为统治者。\BibleQuote{愚人吃饱时：}即轻易得地为业，吃其果实，饱足后（百姓）便踢踹。\BibleQuote{婢女赶走主母时：}即夺取主生命、将基督的肉体钉在十字架上的会堂。\par

\BibleQuote{地上有四样最小之物，却比智者更智慧：蚂蚁没有力量，却在夏天预备粮食。}同样，外邦人藉着对基督的信德，藉善功为自己预备永生。\BibleQuote{石獾是软弱之民，却在岩石中造房屋。}也就是说，外邦人建立在基督——这成为房角石的属神磐石——之上。\BibleQuote{蜘蛛用手支撑自己，轻易被捉拿，却住在君王的堡垒中。}即那（在十字架上）伸开双手的强盗，安息于基督的十字架上，住在三位君王——圣父、圣子、圣神——的堡垒，即乐园中。\par

\BibleQuote{蝗虫没有君王，却如受一人指挥般列队前进。}外邦人没有君王，因他们被罪辖制；但现在，他们信从天主，从事天上的争战。\par

\BibleQuote{有三样行走优美，第四样行走优雅：}即天上的天使、地上的圣人、地下的义人灵魂。第四样，即降生成人的天主圣言，光荣地经过童贞女的胎；祂更新我们的亚当，经过天国之门，成为众人复活与升天的初果。\par

\BibleQuote{狮子的幼崽比野兽更强：}即基督，如雅各伯藉犹大所预言的。\BibleQuote{公鸡在母鸡中昂首行走：}即保禄，在教会中大胆宣讲天主基督的话。\BibleQuote{公山羊带领群羊：}即为世界之罪而献祭的那一位。\BibleQuote{君王在民众中发言：}如此，基督治理万民，并藉先知与宗徒宣讲真理之言。\par

那是固守于邪恶的人。宗徒也说：\BibleQuote{犯罪的人，要在众人面前斥责他们；}即除了被弃绝者之外。\BibleQuote{石獾}所指的，难道不是我们吗？我们从前如同猪，行走在世界的一切污秽中；但现在，我们信从基督，将房屋建造在基督的圣洁肉体上，如同建造在磐石上？\par

大地的震动表示地上事物的转变。因此，本身为奴的罪，曾在人必朽的身体中作王：一次是在洪水时期；另一次在索多玛人的时代，他们不满足于土地所出产的，竟对异乡人行强暴；第三次是在可憎的埃及，虽然它藉若瑟获得了一个分发食物给众人、使众人免于饥荒的人，却未善待他的昌盛，反而迫害以色列子民。\BibleQuote{婢女赶走主母：}即外邦人的教会，她本身为奴，与应许无份，却赶走了生而自由的、作主母的会堂，成为基督的妻子与新娘。藉圣父、圣子、圣神，全地都被震动。\BibleQuote{第四样它不能承受：}因为祂初次藉立法者而来，第二次藉先知，第三次藉福音，公开显现自己；第四次，祂将作为审判生者死者的审判者而来，祂的荣耀，整个受造界都不能承受。\par

% structure: p0133 source=h3 target=subsection reason=relative-heading-rank; base=h2

\subsection{第二片断，论\BibleRef{箴 9:1}}

圣希波律陀论\BibleRef{箴 9:1}，\BibleQuote{智慧建造了房屋。}\par

祂意指基督，即天主圣父的智慧与能力，已建造了祂的房屋，即祂在\Person{童贞玛利亚}所取的肉身本性，正如祂（\Person{若望}）先前所说的：\BibleQuote{圣言成了血肉，并居住在我们中间。}同样，智慧的先知也作证：在世界之前就已存在、生命之源、无限的\BibleQuote{天主的智慧，建造了祂的房屋}——藉着一位不识男子的母亲，即祂取了身体圣殿。\BibleQuote{并立了七根柱子；}即至圣圣神的馨香恩宠，如\Person{依撒意亚}所说：\BibleQuote{天主的七神将安息在祂身上。}但另有人说，七根柱子是指七个神圣的品秩，它们藉着祂圣洁而受默感的教导维系受造界：即先知、宗徒、殉道者、神长、隐修士、圣人与义人。而\BibleQuote{祂宰杀了牲畜，}意指先知和殉道者，他们每日在各城各地为真理被不信者如同羊一般宰杀，并大声呼喊：\BibleQuote{为了祢，我们终日被杀害，被看作待宰的羊群。}再者，\BibleQuote{祂调和了酒}在碗中，意指救主将祂的天主性如纯酒般与\Person{童贞玛利亚}的肉身结合，由她而生，同时是天主与人，二者互不混淆。\BibleQuote{并且摆设了筵席：}这预示了所应许的对天主圣三的知识；也指祂尊贵无玷的身体和血，日日藉着属神的神圣祭台，以祭献的方式施予和奉献，作为那首次且永远可纪念的神圣晚餐之纪念。再者，\BibleQuote{祂派出祂的仆人：}即智慧这样做——也就是基督——以崇高的宣告召唤他们。\BibleQuote{谁若单纯，让他转向我，}她说，显然指圣宗徒们，他们走遍世界，以崇高的神圣宣讲真实地呼召万民认识祂。再者，\BibleQuote{并对缺乏智慧的人说}——即对那些尚未获得圣神能力的人——\BibleQuote{来，吃我的饼，喝我为你调和的酒；}意指祂将祂的神圣血肉与尊贵的血赐给我们，为赦免我们的罪而吃喝。\par

% structure: p0136 source=h2 target=section reason=relative-heading-rank; base=h2

\section{论 \WorkTitle{雅歌}}

1. 北风啊，起来，南风啊，你来，吹在我的花园里，使其中的香料飘散出来（\BibleRef{歌 4:16}）。\Person{若瑟}因这些香料而喜乐，被天主立为王的儿子；\Person{童贞玛利亚}以之受膏，便怀孕了圣言：于是新的奥秘、新的真理、新的国度，以及伟大而不可言说的奥迹都彰显出来了。\par

2. 这一切丰富的知识在哪里？这些奥迹在哪里？这些书在哪里？因为现今仅存的只有《箴言》、《智慧篇》、《训道篇》和《雅歌》。那么，圣经在说谎吗？绝对不可。但他的著作内容多样，正如\Quote{歌中之歌}这句话所表明的；因为那表示他在这本书中汇集了那五千首诗歌的内容。此外，在\Person{希则克雅}时代，有些书被选为用，有些则被搁置。因此圣经说：\BibleQuote{这些是撒罗满的杂集箴言，是希则克雅王的友人抄录的。}他们从哪里抄录呢？岂不是从那包含三千比喻和五千诗歌的书中？于是，\Person{希则克雅}的智慧友人从中选取了那些有益于教会建设的部分。而撒罗满关于\Quote{比喻}和\Quote{诗歌}的书——他在其中论述了植物的生理，以及陆地、空中、海中各种动物的生理，还有疾病的治疗——\Person{希则克雅}却废弃了，因为百姓因这些而寻求疾病的疗方，忽略了从天主寻求医治。\par

% structure: p0139 source=h2 target=section reason=relative-heading-rank; base=h2

\section{论 \WorkTitle{依撒意亚}}

% structure: p0140 source=h3 target=subsection reason=relative-heading-rank; base=h2

\subsection{第一残篇}

罗马主教\Person{依玻理}论\Person{希则克雅}。\par

当犹大王\Person{希则克雅}仍患病痛哭时，有一位天使来对他说：\BibleQuote{我已看见你的眼泪，听见你的声音。看，我要在你的寿数上增加十五年。上主给你一个记号：看，我要使你父亲家台阶上的日影后退，即已下沉的太阳所照的十级，日影后退十级。}使那一天成为三十二小时的一日。原来，太阳已运行到第十小时，却又返回了。此外，当农的儿子\Person{若苏厄}与阿摩黎人交战时，太阳已偏西，战事正紧，\Person{若苏厄}担心夜间来临使异民逃脱，便呼喊说：\BibleQuote{太阳，停在\Place{基贝红}；月亮，停在\Place{阿雅隆谷}！}直到我击败这民族。于是太阳和月亮停在原地，那一日成了二十四小时的一日。在\Person{希则克雅}时代，月亮也与太阳一同回转，以免两个元素天体因互撞而失序。那时，\Place{巴比伦}加色丁王\Person{默洛达客}大为惊异——因为他研究星象学，仔细测量这些天体的运行——得知原因后，便送书信和礼物给\Person{希则克雅}，正如后来东方贤士对基督所做的那样。\par

% structure: p0143 source=h3 target=subsection reason=relative-heading-rank; base=h2

\subsection{第二残篇}

摘自圣\Person{依玻理}论\WorkTitle{依撒意亚}开端的讲道。\par

祂以\Place{埃及}指代世界，以手造之物指代其偶像崇拜，以震动指代其倾覆与瓦解。而主——圣言——被描述为坐在一朵轻云上，意指那最纯洁的帐幕，在其中设立祂的宝座，我们的主耶稣基督来到世界震动谬误。\par

% structure: p0146 source=h3 target=subsection reason=relative-heading-rank; base=h2

\subsection{第三残篇}

我们在前人的注释中发现，那日有三十两个小时。因为太阳行完其程，到了第十小时，殿内的影子已下降了十度，主的话使太阳又往回走了十度，这样就有了二十小时。然后，太阳按通常的规律完成了自己的行程，到达了日落。于是，总共有了三十二小时。\par

% structure: p0148 source=h2 target=section reason=relative-heading-rank; base=h2

\section{论耶肋米亚和厄则克耳}

那么，撒罗满圣殿的尺寸是多少？其长度是六十肘，宽度是二十肘。它不朝向东方，以免朝拜者朝拜升起的太阳，而是朝拜太阳的主。如果有人因圣经说长度是四十肘而我却说是六十肘而感到奇怪，请不要惊讶。因为稍后经文提到了另外二十肘，在描述至圣所时，它也称其为\OriginalTerm{达彼尔}{Dabir}。这样，圣所是四十肘，至圣所又是二十肘。\Person{若瑟夫}说圣殿有两层，总高度是一百二十肘。因为\WorkTitle{编年纪}也这样指示，说：\Quote{撒罗满开始建造天主的殿。起初的长度是六十肘，宽度二十肘，高度一百二十肘；里面贴了纯金。}\par

% structure: p0150 source=h2 target=section reason=relative-heading-rank; base=h2

\section{论达尼尔}

% structure: p0151 source=h3 target=subsection reason=relative-heading-rank; base=h2

\subsection{第一段（序言）}

由于我想准确记述以色列子民在巴比伦被掳的时期，并讨论蒙福的达尼尔在异象中所包含的预言，以及他从幼年在巴比伦的生活方式，我也将为此圣洁、正义的人作证，他是基督的先知和见证，不仅在当时宣告了拿步高王的异象，还培养了与自己志同道合的青年，在世上兴起了忠信的见证。他诞生于蒙福的耶肋米亚先知事奉时期，约雅金（即厄里雅金）为王之时。他与其余被掳的人一同被俘，押往巴比伦。蒙福的约史雅生了以下五个儿子：约阿哈次、厄里雅金、约哈南、漆德克雅（即耶苛尼雅）和撒杜姆。父亲去世后，百姓膏立约阿哈次为王，时年二十三岁。在他登基第三个月，法郎乃苛上来攻打他，将他掳往埃及，并向那地征收一百塔冷通银子、十塔冷通金子作为贡赋。随后，法老立其弟厄里雅金为王治理那地，给他改名为约雅金，当时他十一岁。巴比伦王拿步高上来攻打他，将他掳往巴比伦，同时也带走了耶路撒冷殿中的一些器皿。约雅金因被视为法老的朋友、且由法老立为国王，被投入监牢，最终在第三十七年，巴比伦王厄威耳默洛达客释放了他；他剃了短发，做了王的谋士，与王同席，直到去世之日。约雅金被废后，其子约雅金作王三年。巴比伦王拿步高上来攻打他，将他及其百姓一万人掳往巴比伦，另立他父亲的兄弟为王，给他改名为漆德克雅，并与他起誓立约后，返回巴比伦。漆德克雅作王十一年后背叛，投靠了埃及王法老。第十年，拿步高从加色丁地来攻打他，用栅栏围住城，四面环绕，完全封锁。因此，城中大部分人死于饥荒，另一些人死于刀剑，还有一些人被俘；城被火焚烧，圣殿和城墙被毁。加色丁军队夺取了上主殿中所有财宝，以及一切金银器皿；所有铜器，护卫长乃步匝辣当都剥去，带往巴比伦。加色丁军队追赶趁夜与七百人一同逃跑的漆德克雅，在耶里哥追上了他，把他带到黎贝拉的巴比伦王那里。王因他违背了上主的誓言和自己与他所立的约，发怒审判他；在他眼前杀了他的儿子，并剜了漆德克雅的眼睛。又用铜链锁住他，带往巴比伦；他在那里推磨直到去世。他死后，人将他的尸体抛在尼尼微城墙外。在他身上应验了耶肋米亚的预言：\Quote{上主说：我指着我的永生起誓，犹大王约雅金的儿子耶苛尼雅，虽是我右手上的印玺，我也要将你摘下，交在寻索你性命、你所惧怕的人手中，就是加色丁人手中。我必将你和生你的母亲驱逐到你们不是出生之地，你们必死在那里。但他们心中向往的故土，我必不让你们回去。耶苛尼雅像一件无用的器皿，被人丢弃，被驱逐到不认识之地。大地啊，听上主的话：要写明这人是个被开除的人，因为他的后裔中没有一人能兴旺，坐在达味的宝座上，再在犹大执政。}于是，在出离埃及之后，巴比伦之囚便临到了他们。当全体百姓被掳，城市荒凉，圣所被毁，为要成就上主藉耶肋米亚先知的口所说的话：\Quote{圣所必荒凉七十年；}那时我们发现，蒙福的达尼尔在巴比伦说预言，并作为苏撒纳的辩护者出现。\par

% structure: p0153 source=h3 target=subsection reason=relative-heading-rank; base=h2

\subsection{第二段（论异象）}

罗马主教希波律特对达尼尔和拿步高异象的释义，综合而论。\par

1. 当论及\Quote{从海里上来的母狮}时，他指的是\Place{巴比伦}王国的兴起，这也是像的\Quote{金头}。论及它的\Quote{鹰翅}时，他意指\Person{拿步高}王被高举，他的荣耀被提升以对抗天主。然后他说\Quote{它的翅膀被拔去}，即他的荣耀被摧毁；因为他被逐出了他的王国。而\Quote{兽心被换成人心，并叫它用两脚站立，像人一样}这句话，意思是它又恢复了理智，认识到自己不过是个人，并将荣耀归给天主。此后，在母狮之后，他看见第二只兽\Quote{如同熊}，这象征了\Place{波斯}人。因为\Place{巴比伦}之后，\Place{波斯}人获得了权力。而说\Quote{它口中衔着三根肋骨}时，他指的是三个民族：\Place{波斯}人、\Place{玛待}人和\Place{巴比伦}人，这些在像中以金后的银子来表示。接着是第三只兽\Quote{豹子}，意指\Place{希腊}人；因为\Place{波斯}之后，\Place{马其顿}的\Person{亚历山大}获得了权力，当时\Person{达理阿}被推翻，这也由像中的铜所表明。而说那兽\Quote{有四个鸟翅和四个头}时，他非常清楚地表明\Person{亚历山大}的王国分裂为四部分。因为提到四个头时，他指的是从中兴起的四个王。原来\Person{亚历山大}在临终时，将他的王国分成了四部分。然后他说：\Quote{第四只兽是可怖而可怕的；它有铁牙和铜爪。}这指的是谁呢？不就是\Place{罗马}人吗？他们的王国——至今仍屹立的王国——由铁所代表；因为他说：\Quote{它的腿是铁的。}\par

2. 此后，亲爱的，剩下的不就是那像的脚趾，其中\Quote{一部分是铁，一部分是泥，彼此混杂}吗？他奥秘地以脚趾指那从这王国中兴起的十个王。正如\Person{达尼尔}所说：\Quote{我注视那兽，看，在它后面有十只角，其中要另兴起一只小角，从它们中间生出。}这所指的，正是那将要兴起的假基督；他要建立犹大王国。而说由这一只角\Quote{拔起了三只角}，他指明\Place{埃及}、\Place{利比亚}和\Place{埃塞俄比亚}的三个王，这一只角要在战争的阵列中杀死他们。当他征服一切后，他将证明自己是一个可怕而残暴的暴君，并对圣人们施加苦难和迫害，高举自己反对他们。在他之后，就有\Quote{石头}从天而降，\Quote{击打那像}，将其打碎，推翻所有王国，并将王国交给至高者的圣人们。这石头\Quote{变成一座大山，充满了全地。}\par

3. 既然这些事注定要发生，而像的脚趾成为民主政体，兽的十角分配给了十位王，让我们更仔细地查看眼前的事，仿佛睁着眼睛审视它。像的\Quote{金头}与\Quote{母狮}相同，\Place{巴比伦}人由此被代表。\Quote{金的肩膀和银的手臂}与\Quote{熊}相同，\Place{波斯}人和\Place{玛待}人由此被意指。\Quote{铜的肚腹和大腿}是\Quote{豹子}，意指\Place{希腊}人，即从\Person{亚历山大}开始的统治者。\Quote{铁腿}是\Quote{可怖而可怕的兽}，意指现在统治帝国的\Place{罗马}人。\Quote{泥和铁的脚趾}是将来要出现的\Quote{十角}。\Quote{从它们中间另兴起的一只小角}是\Quote{假基督}。那块\Quote{击打像，将像砸碎}并充满全地的石头，是基督，他从天而来，对世界施行审判。\par

4. 但为免我们的主题在此未得证明，我们必须讨论时间的问题，对此人不应轻率发言，因为时间是他的光。因为时间从世界的奠基起被记录，并从\Person{亚当}算起，它们清晰地向我们呈现了我们的探究所涉及之事。原来我们的主在肉身内的首次显现，发生在\Place{白冷}，在\Person{奥古斯都}年间，年代为5500年；他在第三十三年受难。六千年必须完成，以便安息日到来，即那安息、圣日\Quote{天主在这一天歇了祂所做的一切工程}。因为安息日是未来圣者王国——当他们与基督一同为王时——的预像和记号，那时基督从天降临，正如\Person{若望}在他的\WorkTitle{默示录}中所说：因为\Quote{在主前一日如千年}。既然天主在六日内创造了万物，结论就是六千年必须完成。它们尚未完成，正如\Person{若望}所说：\Quote{五位已经跌倒，一位还在}——即那第六位——\Quote{另一位还没有来到。}\par

5. 此外，在提到\Quote{另一位}时，他特别指出了第七位，即安息所在。但有人或许会问：你如何向我证明救世主生于5500年？人啊，你很容易就能明白；因为古时在旷野中，在\Person{梅瑟}领导下，关于会幕所发生的事，构成了属灵奥秘的预像和象征，以便当真理在末期于基督内来临时，你能明白这些事已应验。因为天主对他说：\Quote{你要用不朽的木制造约柜，内外包上纯金；长二肘半，宽一肘半，高一肘半。}这些尺寸加起来共五肘半，由此便可象征5500年。\par

6. 那时，救主显现，将祂的身体——由童贞玛利亚所生——显示给世界。童贞玛利亚是\Quote{用纯金包裹的约柜}，内有圣言，外有圣神；如此，真理得以显明，\Quote{约柜}得以彰显。因此，从基督诞生起，我们必须算入剩余的\OriginalTerm{五百年}{500 years}，以凑足六千年，然后终结将至。救主在第五个半日（即第五日半）显现于世界，带着不朽的约柜，即祂的身体；正如若望所宣告：\BibleQuote{那时是第六个时辰}\BibleRef{若 19:14}，他说这话暗示白日的一半。但主的一日是一万年；因此其半数是五百年。因为祂不应早些显现，因为法律的负担仍然存留；也不应在第六日完满时显现（因为洗礼已改变）；而是在第五个半日，好使在剩余的一半时间里，福音能传遍全世界，并在第六日完满时，祂终结今世的生命。\par

7. 既然波斯人统治了三百三十年，之后更荣耀的希腊人统治了三百年，那么第四兽既强大又胜过以前的一切，必然要统治五百年。当时期满全，在最后（时期）从兽中生出十角时，\OriginalTerm{假基督}{Antichrist}将在他们中出现。当他向圣徒作战并迫害他们时，那时我们可以期待主从天上显现。\par

8. 先知既已如此精确地教导我们关于将来之事的确定性，便中断了当前的话题，再次转到波斯和希腊的王国，向我们讲述另一个异象，这个异象发生了，并在适当的时候应验了；为的是藉此建立我们的信念，使他能向天主呈献我们，使我们对将来之事更易相信。因此，他在第一个异象中所叙述的，再次详细讲述，以建立信友的德行。因为藉着\BibleQuote{那向西、向北、向南撞的公绵羊}\BibleRef{达 8:4}，他指的是波斯王\Person{达理阿}，他征服了万民；他说：\BibleQuote{这些兽不能站立在他面前}\BibleRef{达 8:7}。而藉着\BibleQuote{那从西方而来的公山羊}\BibleRef{达 8:5}，他指的是希腊王\Person{亚历山大}（马其顿人）；至于它\BibleQuote{来到那公绵羊面前，怒气冲冲，击打它的脸，将它打碎，抛在地上，并践踏它}\BibleRef{达 8:7}，这正是所发生的事。\par

9. 因为\Person{亚历山大}与\Person{达理阿}争战，战胜了他，并在击溃并摧毁他的营寨后，使自己成为整个统治权的主人。然后，在公山羊高升之后，它的角——就是那大角——折断了；在它下面长出四只角，向着天的四方。因为当\Person{亚历山大}控制了所有波斯土地，并使人民臣服后，他随即去世，将他的王国分为四个公国，如上所述。从那以后，\BibleQuote{一只角被高举，壮大到天上的权势；藉着它，祭献受到扰乱，正义被抛到地上}\BibleRef{达 8:10-11}。\par

10. 因为\Person{安提约古}兴起，别号\OriginalTerm{厄丕法乃}{Epiphanes}，是\Person{亚历山大}的后裔。他在叙利亚统治，并使整个埃及臣服后，上到\Place{耶路撒冷}，进入圣所，夺取了上主殿中的所有宝藏、金灯台、桌子、祭坛，并在国内大行杀戮；正如记载：\BibleQuote{圣所要被践踏，直到晚上和早晨，共一千三百天}\BibleRef{达 8:14}。因为圣所确实在那段时期荒凉了三年半，以使一千三百天得以满全；直到\Person{犹大玛加伯}在他父亲\Person{玛塔提雅}死后兴起，抵挡他，摧毁了\Person{安提约古}的营寨，解救城市，收复圣所，并严格按照法律修复圣所。\par

11. 既然天使\Person{加俾额尔}也向先知叙述了这些事，正如我们所理解的，也正如它们已经发生，并都在\WorkTitle{玛加伯书}中清楚记述，让我们进一步看他关于其他周所说的话。因为当他读了\Person{耶肋米亚}先知的书，其中记载圣所要荒凉七十年，他带着禁食和恳求认罪，祈祷人民能早日从被掳中归回\Place{耶路撒冷}城。于是，他在记述中这样说：\BibleQuote{在\Person{玛待人}\Person{达理阿}元年，他是\Person{阿哈随鲁}的儿子，出身\Person{玛待}的后裔，作\Person{加色丁人}国的王；我——\Person{达尼尔}——从书上明白，主的圣言临到\Person{耶肋米亚}先知，论到\Place{耶路撒冷}荒凉的年数，七十年为满}\BibleRef{达 9:1-2}，等等。\par

12. 在他认罪和恳求之后，天使对他说：\BibleQuote{你是极蒙眷爱的人}\BibleRef{达 9:23}；因为你渴望看见的事，我将指示你；这些事将在它们的时候应验；他触摸我，说：\BibleQuote{为你的人民和你的圣城，已经定了七十个七，要封住罪恶，除净过犯，封住异象和先知，并膏抹至圣所；你当知道、当明白，从发出命令回复，并重建\Place{耶路撒冷}，直到基督君王，有七个七和六十二个七}\BibleRef{达 9:24-25}。\par

13. 因此，在提及七十个星期，并将其分为两部分后，为了使他对先知所说的话更易理解，他接着说：\BibleQuote{直到受傅者君王，有七个星期}\BibleRef{达 9:25}，这共计四十九年。第21年是\Person{达尼尔}在\Place{巴比伦}看到这些事的那一年。因此，四十九年加上二十一年，正好是七十年，真福\Person{耶肋米亚}曾论及这七十年说：\Quote{从他们在\Person{拿步高}下所受的充军起，圣所将荒凉七十年；此后，百姓将回归，献上祭品和供物，那时基督是他们的君王。}\par

14. 那么，他所说的基督是谁呢？岂不是\Person{耶叔亚}，\Person{约匝达克}的儿子，他那时与百姓一同返乡，并在第七十年圣殿重建时，按照法律献上祭品？因为所有的君王和司祭都曾被称为受傅者（基督），因为他们是用\Person{梅瑟}古时所制的圣油受傅的。他们以自己的身份承载主的名字，预表了那从天上而来的完美的君王和司祭，只有他独自承行父的旨意；正如\BibleBook{列王}{纪}所载：\BibleQuote{我要为我兴起一位忠信的司祭，他将按照我的内心行事。}\BibleRef{撒上 2:35}\par

15. 因此，为了指明真福\Person{达尼尔}渴望见到的那一位何时来临，他说：\BibleQuote{过了七个星期，还有六十二个星期}\BibleRef{达 9:25}，这一时段包含四百三十四年。因为自百姓在\Person{耶叔亚}（\Person{约匝达克}的儿子）、经师\Person{厄斯德拉}以及达味支派\Person{沙耳}的孙子\Person{则鲁巴贝耳}领导下从\Place{巴比伦}返回，直到基督来临，共有四百三十四年，以便众司祭的司祭能在世间显现，那位除免世罪者能明确彰显，正如\Person{若翰}论及他所说：\BibleQuote{看，天主的羔羊，除免世罪者！}\BibleRef{若 1:29}。同样，\Person{加俾额尔}说：\BibleQuote{为终止过犯，为罪恶做补赎。}\BibleRef{达 9:24}。但谁终止了我们的过犯？宗徒\Person{保禄}教导我们说：\BibleQuote{他是我们的和平，使双方合而为一；}\BibleRef{弗 2:14} 然后：\BibleQuote{涂抹了那反对我们的债卷。}\BibleRef{哥 2:14}\par

16. 因此，过犯被涂抹，罪恶得到补赎，由此得以显明。但谁是为他们的罪得到补赎的人呢？岂非那些信他名，并以善行博取他悦纳的人吗？自百姓从\Place{巴比伦}返回，直至基督诞生，其间有四百三十四年，这很容易理解。因为，既然第一次盟约是在四百三十四年后赐予以色列子民的，那么第二次盟约也应以相同的时间界定，以便百姓期待，并容易被信友认出。\par

17. 为此，\Person{加俾额尔}说：\BibleQuote{并要傅抹至圣所。}\BibleRef{达 9:24}。而至圣所不是别的，正是天主子自己。当他来临并显现时，对他们说：\BibleQuote{上主的神临于我，因为他给我傅了油；}\BibleRef{路 4:18} 等等。因此，凡相信天上司祭的人，都因这同一位司祭而被洁净，他们的罪也被涂抹。凡不相信他的人，轻视他如同凡人，他们的罪就被封存，作为不能除去之物；因此，天使预见到并非所有人都信他，就说：\BibleQuote{为终结罪恶，并为密封罪恶。}\BibleRef{达 9:24}。因为凡继续不信他的人，直到最后，他们的罪不是被终结，而是被密封保存以待审判。但凡相信他能赦罪的人，他们的罪就被涂抹。因此他说：\BibleQuote{并为密封异象和先知。}\BibleRef{达 9:24}\par

18. 因为当那一位——法律和先知的圆满（因为法律和先知直到\Person{若翰}为止）——来临时，由他们所说的一切必须得到确认（密封），以便在主来临时，一切被释放的都得以显明，而古时被捆绑的现由他释放，正如主亲自对百姓的领袖所说，当他们对安息日的治愈感到愤怒时：\BibleQuote{假善人！你们中谁不在安息日解开自己的牛或驴，牵去饮水？这女人是\Person{亚巴郎}的女儿，撒殚捆绑了她十八年，难道不该在安息日解开她的束缚吗？}\BibleRef{路 13:15-16}。因此，凡被撒殚用锁链捆绑的，主来临时就解除了死亡的束缚，捆绑了我们的强敌，释放了人类。正如\Person{依撒意亚}所说：\BibleQuote{然后他对被囚的说：出来！对在黑暗中的说：显露吧！}\BibleRef{依 49:9}\par

19. 古时由法律和先知所说的一切都被密封，不为人所知，\Person{依撒意亚}宣告说：\BibleQuote{他们将那密封的书卷交给识字的，说：读读这个！他回答说：我不能读，因为它是密封的。}\BibleRef{依 29:11}。古时先知所说的一切，对于不信的法利塞人——他们自以为理解法律的文字——是密封的，而对于信者则是开启的，这是合适而必要的。因此，古时密封的事，如今因天主的恩宠，都为圣徒开启了。\par

20. 因为祂自己就是完美的印，教会是钥匙：\Quote{祂开，无人能关；祂关，无人能开，}正如\Person{若望}所说。再者，同一位又说：\Quote{我看见在那坐于宝座者右手上，有一书卷，内外都写着字，用七个印封住；我又看见一位天使大声宣布说：谁配展开那书卷，揭开它的印呢？}等等。\Quote{我又看见在宝座和四个活物中间，站着一只羔羊，好像被杀过的，有七角七眼，这些是天主被差遣往普天下的七神。祂前来从坐于宝座者右手中取了那书卷。当祂取了书卷，四个活物和二十四位长老就俯伏在羔羊面前，各拿着琴和盛满了香的金炉，这香就是圣人们的祈祷。他们唱新歌，说：祢配拿那书卷，揭开它的印，因为祢曾被宰杀，用祢的血赎了我们归于天主。}于是祂取了书卷，并且揭开了，为要使那些从前关于祂隐秘所说的事，如今可以在屋顶上坦然宣布。\par

21. 因此，天使对\Person{达尼尔}说：\Quote{封闭这话，因为这神视要到时期的终结。}但对基督却不是说\Quote{封闭}，而是\Quote{解开}那从前捆绑的事；为使我们借祂的恩宠，能知道圣父的旨意，并相信祂所派遣来为人的救恩的那一位，即耶稣我等主。因此，他说：\Quote{他们必回归，街道必重建，城墙必重建，}这确实发生了。因为百姓回归，建造了城、圣殿和周围的城墙。然后他说：\Quote{过了六十二周，时期将要满全，有一周要与许多人订立盟约；在一周的中段，祭献和奉献将被除去，在圣殿中必有可憎的荒凉。}\par

22. 因为当六十二周满全，基督来到，福音在各处被宣讲，时期那时完成，只余下一周，就是最后一周，其中\Person{厄里亚}和\Person{哈诺客}将要出现，而在其中段，可憎的荒凉将显现出来，即假基督，向世界宣告荒凉。当他来到时，祭献和奉献将被除去，就是如今万邦在各处向天主所献的。这些事既如此叙述，先知又向我们描述另一个神视。因为他没有别的挂虑，只求准确了解所有将发生的事，并证明自己是这些事的导师。\par

23. 于是他说道：\Quote{在波斯王\Person{居鲁士}第三年，有话启示给\Person{达尼尔}，他名叫\Person{贝耳沙匝}；这话是真实的，并有伟大的能力和明悟在神视中赐给他。在那些日子，我\Person{达尼尔}悲痛了三周的日子。我没有吃美味的食物，肉和酒没有入我的口，也没有用油抹我自己，直到三周的日子满。正月第四日，我谦卑了自己，}他说，\Quote{二十一天，}向永生天主祈祷，求祂启示这奥迹。圣父果真俯听了我，并派遣了祂的圣言，藉祂显明将要发生的事。这事确实发生在大河旁。因为圣子在那里显现是相宜的，祂也在那里要除免罪过。\par

24. \Quote{我举目观看，}他说，\Quote{看，有一个人身穿细麻衣。}在第一个神视中，他说：\Quote{看，\Person{加俾额尔}天使被派遣。}但在这里却不然；他看见的是主，虽然尚未成为完人，却具有人的外貌和样式，正如他所说：\Quote{看，有一个人身穿细麻衣。}因为他身穿彩色外衣，奥秘地指示了我们蒙召的多种恩宠。因为司铎袍是由不同颜色制成的，表明万邦等候基督的来临，为使我们能像由许多颜色组成的一个身体。\Quote{他的腰束着\OriginalTerm{纯金}{Ophaz}。}\par

25. 现在\OriginalTerm{纯金}{Ophaz}一词是从希伯来语转到希腊语的词，意指纯金。因此，他腰间束着纯金的腰带。因为圣言要担负我们众人，像腰带一样将我们束在祂的身体上，系于祂自己的爱中。完整的身体是祂的，而我们都是祂身体上的肢体，互相结合，并由圣言自己维系。\Quote{他的身体如同\OriginalTerm{埃塞俄比亚人}{Tharses}。}如今\OriginalTerm{埃塞俄比亚人}{Tharses}，按解释，就是\Quote{埃塞俄比亚人}。因为人们难以认出祂，先知早已预先宣告，暗示祂要在肉身中显现于世界，但许多人会难以认出祂。\Quote{他的面貌如同闪电，眼睛如同火灯；}因为圣言如火和审判的权能理应预先被表明，在行使这权能时，祂将叫（祂审判的）火按公义临到不虔敬者身上，并消灭他们。\par

26. 他还加上这些话：\Quote{他的臂和脚如同磨亮的铜；}以表示人的初次蒙召，和与之相似的第二次蒙召，即外邦人的蒙召。\Quote{因为后来的要如同先前的；因为我要立你的官长如起初，你的首领如以前。祂的声音如同大群众的声音。}因为我们所有在这些日子相信祂的人，都发出预言般的话，如同借祂的口说出祂所指定的事。\par

27. 过了一会儿，祂对他说：\Quote{你知道我为什么到你这里来吗？现在我要回去与\OriginalTerm{波斯的君}{prince of Persia}争战。但我将指示你\OriginalTerm{真理的书}{Scripture of truth}上所记的事：在这些事上，除了你的君米迦勒之外，没有人支持我，我把他留在那里。因为从你决心在主你的天主面前受苦的那一天起，你的祈祷已蒙垂听，我奉差遣来与波斯的君争战。}因为有一个计谋形成，要不释放百姓：\Quote{因此，为使你的祈祷迅速获允，我抵挡了他，并把米迦勒留在那里。}\par

28. 那说话的是谁，岂不就是赐给百姓的天使，正如\OriginalTerm{梅瑟法律}{law of Moses}中所说：\Quote{我不与你们同行，因为百姓是硬颈的；但我的天使要在你们前面同行。}这位天使在旅店中阻挡了梅瑟，当时他正将未受割礼的孩子带往埃及。这是因为梅瑟，作为法律的长者（或使节）和中介，宣告了祖先的\OriginalTerm{盟约}{covenant}，不允许引进未受割礼的孩子，免得被百姓视为假先知和欺骗者。\Quote{如今，}他说，\Quote{我要将真理指示你。}真理岂能指示别的，只能指示真理吗？\par

29. 因此他对他说：\Quote{看，在波斯还要兴起三位君王；第四位要远比其他都富有；当他获得财富后，要起来攻击希腊的一切王国。必有一位勇武的君王兴起，执掌大权，任意而行；当他的王国建立后，必被破碎，向\OriginalTerm{天的四方}{four winds of heaven}分裂。}这些事我们已经在前面论及\OriginalTerm{四兽}{four beasts}时讨论过了。但由于圣经现在再次明确陈述这些事，我们也要第二次讨论它们，免得我们闲置圣经而不加以解释。\par

30. \Quote{还要兴起三位君王，}他说，\Quote{在波斯；第四位要远比其他都富有。}这已经应验了。因为在\Person{居鲁士}之后，兴起\Person{达理阿}，然后是\Person{阿塔薛西斯}。这就是那三位君王；圣经便应验了。\Quote{第四位要远比其他都富有。}那不是别人，正是\Person{达理阿}，他统治并使自己显赫——他富有，并攻击了希腊的一切领域。\Person{马其顿的亚历山大}起来攻击他，毁灭了他的王国；在他征服波斯人之后，他自己的王国就向\OriginalTerm{天的四方}{four winds of heaven}分裂了。因为亚历山大在死时将自己的王国分成了四个公国。\Quote{必有一位君王兴起，进入埃及王的堡垒。}\par

31. 因为\Person{安提约古}成了叙利亚王。他在希腊王国第一百零七年执掌主权。就在那些时期，他与埃及王\Person{仆托肋米}交战，并征服了他，取得了权势。从埃及返回时，他于第一百零三年上耶路撒冷，并带走了主殿的所有财宝，前往\Place{安提约基雅}。过了两年，王派他的税吏前往犹太的各城，强迫犹太人离弃他们祖先的法律，服从王的谕旨。他来并试图强迫他们说：\Quote{出来，遵行王的命令，你们就得以存活。}\par

32. 但他们说：\Quote{我们不出来，也不遵行王的命令；我们愿死在自己的清白中。}他便杀了他们一千人。因此，对蒙福的\Person{达尼尔}所说的话都应验了：\Quote{我的仆人将受苦难，因饥荒、刀剑和掳掠而倒下。}达尼尔却补充说：\Quote{他们必得少许的帮助。}因为那时\Person{玛塔提雅}兴起，还有\Person{犹大玛加伯}，帮助他们，并从希腊人手中拯救了他们。\par

33. 因此，圣经中所说的话应验了。他接着这样说：\Quote{南方王（的女儿）要前往北方王那里，与他立约；那引领她的人的手臂不能站立；她也要被击打，跌倒，那引领她的人也是如此。}因为这是埃及的王后\Person{仆托肋买}。那时，她带着两个儿子\Person{仆托肋米}和\Person{斐洛默托尔}，前去与叙利亚王\Person{安提约古}立约；当她到达\Place{西托波利斯}时，在那里被杀了。因为那引领她的人出卖了她。就在那时，两兄弟彼此争战，\Person{斐洛默托尔}被杀，\Person{仆托肋米}掌握了权力。\par

34. 然后，\Person{仆托肋米}再次向\Person{安提约古}开战，\Person{安提约古}迎战他。因为圣经这样说：\Quote{南方王要起来攻击北方王，她的后裔也要起来攻击他。}是什么后裔，不就是\Person{仆托肋米}吗？他与\Person{安提约古}交战。\Person{安提约古}出兵攻击他，未能战胜他，只得逃跑，返回\Place{安提约基雅}，并召集了更多的军队。\Person{仆托肋米}便夺取了他的一切装备，运往埃及。圣经便应验了，如达尼尔所说：他要将他们的神像、铸像和一切珍贵的金器运往埃及。\par

35. 此后，\Person{安提约古}第二次出兵攻打他，并战胜了\Person{仆托肋米}。这些事件之后，\Person{安提约古}再次向以色列子民开战，派遣一个名叫\Person{尼加诺尔}的人率领大军去征服犹太人。那时，\Person{玛塔提雅}死后，\Person{犹大玛加伯}统治着百姓；其余的事，记载在\WorkTitle{玛加伯书}上。这些事发生后，圣经又说：\Quote{必有一位君王兴起，在地上得胜；南方王要兴起，并娶他的女儿为妻。}\par

36. 因为曾发生这样的事：有一个名叫\Person{亚历山大}的人，是\Person{斐理伯}的儿子。他当时对抗\Person{安提约古}，与他作战，并击败了他，夺取了王国。然后他派人去见埃及王\Person{仆托肋米}，说：\Quote{把你的女儿\Person{克娄巴特拉}给我为妻。}他便将\Person{克娄巴特拉}嫁给\Person{亚历山大}为妻。这样，圣经便应验了，它说：\BibleQuote{他必得他的女儿为妻。}圣经进一步说：\BibleQuote{他必玷污她，她却不会是他的妻子。}这也确实应验了。因为\Person{仆托肋米}将女儿给了他之后，他回去，看见\Person{亚历山大}强大而光荣的王国。他贪图其领土，便对\Person{亚历山大}说谎，正如圣经所说：\BibleQuote{两个王必在（同一）席上说谎。}事实上，\Person{仆托肋米}回到埃及，集结大军，在\Person{亚历山大}进军\Place{基里基雅}时攻打那城。\par

37. \Person{仆托肋米}于是入侵该地，在各城设立驻军；在占领犹太后，他就出发去见他的女儿，并派人送信给在岛屿上的\Person{德默特琉}，说：\Quote{你到这里来见我，我将把我的女儿\Person{克娄巴特拉}给你为妻，因为\Person{亚历山大}想杀我。}\Person{德默特琉}便来了，\Person{仆托肋米}接待了他，将原本许配给\Person{亚历山大}的女子给了他。这就应验了所记载的话：\BibleQuote{他必玷污她，她却不会是他的妻子。}\Person{亚历山大}被杀。然后\Person{仆托肋米}戴上了两顶王冠，即叙利亚和埃及的王冠，并在戴上后第三天去世。这就应验了圣经所记载的话：\BibleQuote{他们必不将王国的荣耀给他。}因为他死了，没有从众人那里得到作王的尊荣。\par

38. 先知在这样讲述了已经发生并在它们的时候应验了的事情之后，又向我们宣告另一奥秘，同时指出末后的时期。因为他这样说：\BibleQuote{必另有一个无耻的王兴起；他必高举自己超过一切神明，自大，说奇事，并亨通直到怒气成就；}等等。\BibleQuote{这些必从他的手中逃脱：\Place{厄东}、\Place{摩阿布}和\Place{阿孟}子民的首领（或权柄）。他必伸手攻击那地，\Place{埃及}地也不能逃脱。他必掌管金银的隐秘宝藏，以及\Place{埃及}和\Place{利比亚}人及\Place{埃塞俄比亚}人在其堡垒中的一切珍宝。}\par

39. 因此，先知这样论及\OriginalTerm{假基督}{Antichrist}，他将是无耻的，好战的，专制的，他高举自己超过一切君王和一切神明，要建造\Place{耶路撒冷}城，并修复圣所。不敬神的人将敬拜他为神，向他屈膝，以为他就是基督。他将剪除基督的两个见证人和前驱，他们宣告基督从天上而来的荣耀国度，正如经上所说：\BibleQuote{我要赐（权能）给我的两个见证人，他们穿着麻衣，预言一千二百六十天。}也正如向\Person{达尼尔}宣告的：\BibleQuote{一个七要与许多人坚定盟约；在七之半，祭献和供物必停止}——这是要表明那一七被分为两半。这两个见证人将宣讲三年半；而\OriginalTerm{假基督}{Antichrist}将在那一七的试炼中与圣徒争战，并使世界荒凉，以便应验所记载的话：\BibleQuote{他们必使那荒凉的可憎之物存留一千二百九十天。}\par

40. 因此，\Person{达尼尔}提到了两个可憎之物：一个是毁灭的，另一个是荒凉的。什么是毁灭的可憎之物呢？岂不就是\Person{安提约古}当时在那里设立的吗？什么是荒凉的可憎之物呢？岂不就是当\OriginalTerm{假基督}{Antichrist}来临时将要普遍发生的事吗？\BibleQuote{这些必从他的手中逃脱：\Place{厄东}、\Place{摩阿布}和\Place{阿孟}子民的首领。}因为他们因着亲属关系与他结盟，并首先拥戴他为王。\Place{厄东}人是以扫的后裔，住在\Place{色依尔}山。\Place{摩阿布}和\Place{阿孟}是从他的两个女儿所出，正如\Person{依撒意亚}也说：\BibleQuote{他们必在外邦人的船上飞翔（扩展自己），也要劫掠海洋；东、西、北的人都必尊敬他们，\Place{阿孟}的子民必首先服从他们。}他将被他们拥立为王，被众人尊大，并成为世界荒凉的可憎之物，统治一千二百九十天。\BibleQuote{那等候而达到一千三百三十五天的是有福的；}因为当那可憎之物来与圣徒作战时，凡存活过他的日子并达到那四十五天，同时另外五十天的时期向前推进的人，天主的国就临到他。\OriginalTerm{假基督}{Antichrist}确实进入那五十天的一部分，但圣徒将同基督继承国度。\par

41. 这些事叙述之后，\Person{达尼尔}继续说：\BibleQuote{看，有两个人站着，一个在河岸这边，一个在河岸那边；他们回答站在河岸上的那人说：\InnerQuote{这些奇妙的话要到何时才结束呢？}我听见那身穿细麻衣、站在河水之上的人，向天举起右手和左手，指着永生者起誓说：\InnerQuote{要经过一段时期、两段时期和半段时期；当分散完成时，他们必明白这一切事情。}}\par

42. 那么，站在河岸上的两个人是谁呢？不就是法律和先知吗？而站在水上的那位是谁呢？不就是他们古时所预言的那位，在末世在约但河被圣父所见证，并被若翰放胆向百姓宣告的那位吗？祂\Quote{腰间佩带文士的\OriginalTerm{casty}{casty}，身穿杂色细麻衣}。因此，他们向祂询问，知道祂已被赋予一切统治和权能，为要精确地从祂得知，祂何时将审判带到世上，以及祂所说的话何时应验。祂为了无论如何要说服他们，便举起右手和左手向天，指着那永生者起誓。谁在起誓？又指着谁起誓？显然是子指着父起誓，说：\Quote{父永生，但在一个时期，两个时期，半个时期，当分散完成时，他们必要知道这一切。}\par

43. 祂伸出双手，预示了祂的苦难；而提到\BibleQuote{一个时期，两个时期，半个时期，当分散完成时}，祂指出了假基督的三年半。因为祂用\Quote{一个时期}指一年，\Quote{两个时期}指两年，\Quote{半个时期}指半年。这就是达尼尔所预言的\OriginalTerm{一千二百九十日}{thousand two hundred and ninety days}，为完成苦难，并在假基督来临时完成分散。在那日子，他们必要知道这一切。从日常祭被废除的时候算起，也有一千二百九十日。那时罪恶满盈，正如主也说：\BibleQuote{因为罪恶满盈，许多人的爱德就冷淡了。}\par

44. 分裂必然在背教发生时出现。当分裂出现时，爱德便冷淡了。\BibleQuote{那等候并来到一千三百三十五日的人是有福的}这句话也有其意义，正如主所说：\BibleQuote{那坚持到底的，必然得救。}因此，我们绝不可接受背教，免得罪恶满盈，荒凉的可憎之物——即那仇敌——追上我们。祂对他说：\BibleQuote{直到晚上}——即直到终结和早晨。什么是\Quote{早晨}？复活的日子。因为那是另一个时代的开始，正如早晨是一天的开始。而一千四百天是世界的光。因为当光出现在世界时（如祂说：\BibleQuote{我是世界的光}），圣所将被洁净，正如祂所说，脱离那仇敌。因为除非消灭那仇敌，圣所绝不能洁净。\par

% structure: p0199 source=h3 target=subsection reason=relative-heading-rank; base=h2

\subsection{第三片段（达尼尔书注释）}

\BibleRef{达 1:1} \BibleQuote{在约雅金王在位第三年。}圣经叙述这些事，为要暗示百姓第二次被掳，那时约雅金、与他同行的三个青年以及达尼尔一同被掳走。\par

2. \BibleQuote{主将……交在……手中}等。这些话，\BibleQuote{主将……交在……手中}被记录下来，为的是使人读这卷书开篇时，不将他们的被掳归因于掳掠者的强大和首领的懈怠。此外，\BibleQuote{一部分}说得很好，因为被掳是为纠正，而非摧毁整个民族，免得人们对原因产生误解。\par

8. \BibleQuote{达尼尔心中决意。}哦，有福的是那些如此遵守祖传盟约、不违犯梅瑟所传的法律、而是敬畏那位藉他启示的天主的人。这些人虽在异乡为俘虏，却不被美食所诱，也不作酒乐的奴隶，更不被王侯荣耀的诱饵所捕获。他们保守自己的口圣洁纯净，使纯洁的言语从纯洁的口中发出，并用这样的口赞美天上的圣父。\par

12. \BibleQuote{请现在就试验你的仆人。}他们教导说，并非地上的食物使人美丽强壮，而是由圣言所赐的天主恩宠。\Quote{过了不多时。}你已看见青年们不朽的信德和不变的敬畏之心。他们请求十天的间隔，为证明人若不藉相信主所传的圣言，便不能得天主恩宠。\par

19. \BibleQuote{在他们众人中，没有一人像达尼尔。}这些在巴比伦被证明为忠信的证人，在一切智慧上由圣言带领，为藉他们使巴比伦的偶像蒙羞，拿步高被三个青年所胜，炉火因他们的信德而被遏止，邪恶的长老（或首领）的欲望落空。\par

\BibleRef{达 2:3} \BibleQuote{我作了一个梦。}王所见的梦并非属地的梦，以致能被世上的智者解释；而是属天的梦，按照上智安排和天主的预知，在其适当的时候应验。为此，它被隐藏于思想世俗之事的人，为叫寻求天上之事的人，天上的奥秘得以启示。事实上，在埃及当法郎和若瑟的时候，也有类似的事。\par

5. \BibleQuote{那事已从我遗忘。}为此，异象向王隐藏，为叫天主所拣选的人，即达尼尔，显明为一位先知。因为当一些事对某些人隐藏，而由另一个人揭示时，那讲述的人必然显为一位先知。\par

10. \BibleQuote{他们却说：没有人。}因此，当他们宣告王所询问的事不可能由人说出时，天主指示他们：在人不可能的事，在天主却是可能的。\par

14. \BibleQuote{阿黎约客，王的卫队长}（字面意思是\Quote{大屠夫或厨师}）。因为正如厨师宰杀一切动物并烹调它们，他的职务也有类似的性质。世界的统治者杀害人，如同屠宰野兽一般。\par

23. \Quote{因为祢赐给了我智慧和能力。} 因此，我们应当留意天主的良善，祂如何立刻向配得的人和敬畏祂的人启示并显示（自己），满全他们的祈祷和恳求，正如先知所说：\Quote{谁明智，就会明了这些事？谁聪敏，就会认识它们？}\par

27. \Quote{智者和术士不能。} 他训诫君王不要向世俗之人寻求天上奥秘的解释，因为这些奥秘必由天主在适当的时候完成。\par

29. \Quote{至于你，大王，你的思想。} 因为大王在征服埃及地、占据犹太地区并掳掠百姓之后，在床上思想此后要发生的事；那知晓一切隐秘、洞察人心思的主，藉着像将未来之事启示给他。祂又向他隐藏了异象，为使天主的计划不借巴比伦的智者来解释，而是借蒙福的\Person{达内尔}，作为天主的先知，使隐藏于众人的事得以显明。\par

31. \Quote{看，有一座巨大的像。} 那么，我们怎能不留意\Person{达内尔}在巴比伦所预言、如今仍在世上逐步应验的事呢？当时向\Person{拿步高}显示的像，正是全世界的预像。那时，巴比伦人统治万有，他们就是像的金头。其后，波斯人统治了二百四十五年，由银代表。接着，希腊人从马其顿的\Person{亚历山大}开始统治了三百年，因此他们是铜。此后，罗马人到来，他们是像的铁腿，因为他们坚如铁。然后有泥和铁的脚趾，象征着随后兴起的民主政体，分散在像的十个脚趾中，其中铁与泥混合。\par

31. \Quote{你看见，}等等。\Person{阿波利纳尔}论及此事：他观看，看，仿佛一座像。因为那并非作为实际物体呈现在旁观者眼前，而是一个像或幻像。它虽包含许多事物于一身，却并非真正的一体，而是多元的。因为它囊括了所有王国的概要；其极度光辉是因君王的荣耀，其恐怖外表是因他们的权势。\Person{欧瑟伯·旁斐利}和至圣的罗马主教\Person{依玻理}将这里所论的\Person{拿步高}的梦与先知\Person{达内尔}的异象相比较。由于他们在各自的注释中对眼前的异象给出了不同的解释，我认为有必要转录凯撒勒雅的\Person{欧瑟伯}（别号旁斐利）在其\WorkTitle{福音证明}第十五卷中的论述；他以如下方式解释整个异象：我认为此异象（即\Person{拿步高}的梦）与先知的异象毫无差别。因为正如先知看见一片大海，大王也看见一座巨大的像。同样，先知看见四只兽，解释为四个王国，大王也藉着金、银、铜、铁明白四个王国。再者，正如先知看见末后兽的十个角被一只角折断三个，大王也同样看见像的末后部分一部分是铁，一部分是泥。此外，正如先知在看见四个王国的异象之后，看见人子获得统治权、能力和王国，大王也看见一块石头击碎整个像，变成一座大山充满全地。这是恰当的。因为大王对人生景象的看法如此虚妄，他欣赏那幅画中仅感官色彩的美丽，好像把它比作一座巨大的像；而先知则把人生巨大而汹涌的骚动比作汹涌的大海。同样，大王珍视世人视为宝贵的物质（金、银、铜、铁），便将不同时期统治人世的王国比作这些物质；而先知则按照这些王国统治的方式，将它们描述为兽的样式。再者，大王——他似乎骄傲自大，以祖先的权势自夸——看到了世事变迁和地上所有王国注定的结局，借此教导他放下骄傲，明白人间没有稳固之物，只有万物的定终——天主的王国。因为亚述人的第一个王国（以金为表）之后，将有波斯人的第二个王国（以银为表），然后是马其顿人的第三个王国（以铜为表），之后罗马人的第四个王国将接续，比先前的一切更为强大，因此被比作铁。论及此国，经上说：\BibleQuote{第四个王国将坚如铁；如同铁粉碎和征服万物，它也要粉碎和征服万物。} 在这些提及的王国之后，天主的王国由击碎整个像的石头代表。先知与此一致，他在依次描述四兽所象征的四个统治权之前，并未看见在这一切之后到来的王国。我认为，向大王和先知所显示的异象单单关乎这四个王国，而非其他，因为从先知的时代起，犹太民族就受这些王国的奴役。\par

33. \Quote{他的脚，}等等。\Person{依玻理}：在先知的异象中，十个角是将来要发生的事。\par

\\Quote{你观看，直到一块石头被凿出。} 你看见一块仿佛非人手所凿的石头，击打那像的脚。因为人的国度与天主的国度被决定性地分开；关于这事，经上写道：\\Quote{仿佛被凿出。} 然而，那击打击中了末端，并在这些末端上，它摧毁了地上一切统治权。\par

\\Quote{这梦是确定的。} 因此，为使无人怀疑所宣告的事是否会如此发生，先知用这些话证实了它们：\\Quote{这梦是确定的，它的解释是可靠的；} 我并未在解释异象时出错。\par

\\Quote{那时拿步高王俯伏在地。} 拿步高听到这些事，并被提醒他的异象，就知道达尼尔所说的是真的。亲爱的，天主的恩宠能力何其大！那不久前还与其他巴比伦智者一同被判死刑的人，如今竟被王敬拜，不是作为人，而是作为天主！\\Quote{他下令向他献\\OriginalTerm{玛纳}{manaa}}（即，加色丁语称\\Quote{祭品}）\\Quote{和馨香。} 古时，主也曾对梅瑟说过类似的话：\\Quote{看，我使你成为法郎的天主；} 为使梅瑟因在埃及地所行的神迹，不再被视为人，而被埃及人当作天主来敬拜。\par

\\Quote{于是王使达尼尔成为大人物。} 因为他谦卑自己，自视为众人中最卑微的，天主使他伟大，王又立他统管整个巴比伦地。正如法郎对若瑟所做的一样，立他统管全埃及地。\par

\\Quote{达尼尔请求，} 等等。因为他们曾与达尼尔一同向天主祈祷，求祂启示异象给他；所以达尼尔从王获得大尊荣时，也提及他们，向王说明他们所做的，好使他们也因同为见异象者和敬拜天主的人而被视为配受某些尊荣。因为当他们向主求天上的事时，也由王领受了地上的事。\par

\\BibleRef{达 3:1} \\Quote{在第十八年，} 等等。（这些话在拉丁通行本等中缺失。）因此，相当长的时间过去，第十八年正在运行，王想起他的异象，\\Quote{造了一座金像，高六十肘，宽六肘。} 因为当蒙福的达尼尔解释异象时回答王说：\\Quote{你就是那像的黄金头，} 王被这番话冲昏头脑，心中高傲，便复制了这像，好使他自己被所有人当作天主来敬拜。\par

\\Quote{所有民众都俯伏。} 有些人（如此做）是因为惧怕王本人；但所有人（或\\Quote{大多数人}）因为拜偶像，都顺从了王所命令的话。\par

\\Quote{沙得辣客、默沙客、阿贝得乃哥回答说：} 等等。这三个青年给所有忠信之人树立了榜样，因为他们不惧怕总督们的群体，也不因听见王的话而战栗，也不因看见火窑的火焰而退缩，反倒视众人和整个世界为无物，只将敬畏天主摆在眼前。达尼尔虽站在远处保持沉默，却微笑着鼓励他们振作精神。他也为他们所作的见证而喜乐，因他知道这三个青年将在战胜魔鬼中领受荣冠。\par

\\Quote{又下令将窑烧热七倍。} 他下令将大窑烧热七倍，好像他已被他们胜过。就属世的事而言，王更高超；但在对天主的信德上，这三个青年更高超。告诉我，拿步高，你为何下令将他们捆绑着投入火中？难道怕他们脚未被捆绑会逃跑，从而能熄灭火焰吗？但这事并非你自行而为，而是另有其藉着你行这些事。\par

\\Quote{火焰喷射出来。} 他说，火焰从内部被天使驱出，向外迸发。看，连火焰也显得有灵性，仿佛它认出并惩罚了有罪的人。因它不触碰天主的仆人，却烧灭了不信和邪恶的加色丁人。里面的人被天使洒以（清凉的）露水，而那些自以为安全站在窑外的人却被火烧灭。投掷青年的人被火焰烧死，火焰从四面将他们抓住，我想，当他们去捆绑青年时便是如此。\par

\\Quote{第四位形貌似天主子。} 告诉我，拿步高，你何时见过天主子，竟承认这是天主子？谁刺透了你的心，使你说出这样的话？你用什么眼睛能观看这光？为何这事只向你一人显明，而不向你周围的任何总督？但是正如经上所记：\\Quote{王的心在天主手中；} 天主的手在这里，圣言藉此刺透他的心，使他能在窑中认出祂并荣耀祂。我们这意见并非没有充分理由。因为正如以色列子民注定要在世上看见天主却不信祂，圣经预先显示外邦人会认出祂降生成人，而拿步高在古时看见并认出祂未降生成人时在窑中显现，承认祂是天主子。\par

\\Quote{他说道：沙得辣客、默沙客、阿贝得乃哥。} 他这样提名三个青年。但他找不到名字来称呼第四位。因为祂还不是那位由童贞玛利亚所生的耶稣。\par

\\Quote{于是王提拔，} 等等。因为他们因舍身至死而光荣了天主，所以他们自己也受到光荣，不仅受自天主，也受自王。并且他们也教导异族和外邦人敬拜天主。\par

\BibleRef{达 7:1} \BibleQuote{他写下了这梦。}因此，那藉着圣神在异象中启示给这位有福的先知的事，他也都完整地为别人述说了出来，以免他好像只是为自己预言未来，而是为其他也愿意查考神圣圣经的人证明他是先知。\par

2. \BibleQuote{看，有四方的风。}他指受造物的四重划分的存在。\par

3. \BibleQuote{有四只巨兽。}既然当时有各种各样的兽显示给有福的达尼尔，且彼此不同，我们就应当明白，叙述的真实并非涉及某种野兽，而是借着不同野兽的预像和形象，展示了在世界上兴起、统治人类的各种王国。因为他以大海意指整个世界。\par

4. \Quote{直到它的翅膀被拔去。}这事实上发生在\Person{拿步高}时期，正如在前面的书卷中所显示的。他直接见证了这一事在他自己身上应验，因为他被驱逐出王国，被剥去荣誉和他先前所拥有的伟大。\Quote{过了一小会儿：} 那些话：\BibleQuote{它被使如人般用脚站立，并赐予它一颗人心，}这表明\Person{拿步高}，当他谦卑自己，承认自己不过是一个人，服从于天主的权能之下，并向主祈求时，他找到了仁慈，并且被恢复了自己的王国和荣誉。\par

5. \BibleQuote{第二只兽如同熊。}意指波斯王国。\BibleQuote{它有三根肋骨。}他将三个民族称为三根肋骨。因此，意思是：那兽掌管统治，而它之下的这些民族是玛待人、亚述人和巴比伦人。\BibleQuote{有人对它这样说：起来，吞吃。}因为波斯人在那些时期起来，蹂躏了每一块土地，使许多人臣服于他们，并杀害他们。如同这熊是污秽的动物，肉食性，用爪牙撕碎，波斯王国也是如此，它统治了两百三十年。\par

6. \BibleQuote{看，另有一兽如同豹子。}提到豹子，他指希腊人的王国，\Person{马其顿的亚历山大}是他们的王。他将他们比作豹子，因为他们思想敏捷而富有创造性，心里狠毒，如同那动物外表多彩，善于伤害和饮人血。\par

\BibleQuote{这只兽还有四个头。}当\Person{亚历山大}的王国被高举、增长，并在全世界有了名声时，他的王国被分成四个公国。因为\Person{亚历山大}在临终时，将王国分给了同族的四位同伴，即：\Quote{塞琉古、德米特里、托勒密和腓力}；他们都戴上了王冠，正如达尼尔所预言的，也如写在\WorkTitle{玛加伯上}中的。\par

7. \BibleQuote{看，有第四只兽。}如今，在希腊之后，除了现今掌权的之外，没有兴起别的王国，这是众所皆知的。这只兽有铁牙，因为它用力量征服和毁灭一切，如同铁一般。它用脚践踏其余的，因为在这之后没有别的王国，但从它而生出十角。\par

\BibleQuote{它头上有十角。}因为正如先知已经对豹子所说的，那兽有四个头，并且那已经应验，\Person{亚历山大}的王国被分成四个公国，同样现在我们应当期待从它生出的十角，当那兽的时期满全时，那小小的角，即\OriginalTerm{假基督}{Antichrist}，会突然出现在它们中间，正义被逐出大地，整个世界达到它的结局。因此，我们不应当预知天主的计划，而要操练耐心和祈祷，以免落在那样的时期。然而，我们不应拒绝相信这些事将会发生。因为如果先知们先前所预言的事尚未实现，那么我们就不必期待这些事。但如果那些先前的事确实按其所预言的在适当的时候发生了，这些事也必将应验。\par

8. \BibleQuote{我细看这些角。}也就是说，我定睛看着那兽，对有关它的一切感到惊讶，尤其是角的数目。因为这兽的外观在种类上不同于其他兽。\par

13 \BibleQuote{来到万古常存者面前。}藉着\OriginalTerm{万古常存者}{Ancient of days}，他是指主、天主和万有的统治者，以及基督自己，祂使日子变老，自己却不因时间和日子而变老。\par

14. \BibleQuote{他的统治是永远的统治。}圣父将天上地下万有都置于祂自己的圣子权下，藉着一切显示祂为自天主而生的\OriginalTerm{首生者}{first-begotten}，为的是使祂与圣父一同，在天使面前被承认为天主之子，并显现为天使之主；祂也显示为生于童贞女的\OriginalTerm{首生者}{first-begotten}，为叫祂在自己内成为原初受造的亚当的再造者；也显示为从死者中的\OriginalTerm{首生者}{first-begotten}，为叫祂自己成为我们复活的\OriginalTerm{初果}{first-fruits}。\par

\BibleQuote{永不消逝的。}祂展示了圣父赐给祂自己圣子的全部权柄，圣子被显现为天上、地上和地底下万有的君王，以及万有的审判者：在天上的万有，因为祂是\OriginalTerm{圣言}{Word}，在万有之前生于圣父的心；在地上的万有，因为祂成为人，并在自己内重新创造了亚当；在地底下的万有，因为祂也被列于死者之中，向圣人们的灵魂宣讲，并以死亡战胜了死亡。\par

17. \Quote{那将要兴起的。} 因为当三兽结束其行程并被除去，而那一兽仍保持活力时——如果这一兽也被除去，那么最终尘世之事将终结，天上之事开始；为使圣人不可分解、永久的国显现，并让天上的君王向众人彰显，不再是作为图像，如异象中所见，或在山顶云柱中启示的那一位，而是作为天主降生成人和人、天主子和人子，带着天使的权能与军旅，从天降下，作为世界的审判者。\par

19. \Quote{我就询问那第四兽。} 他这里指的是第四个国，我们已经谈过：那个国，在世上没有比它更大的同类国兴起；从它也将生出十角，分配给十个冠冕。在这些角中，另有一小角要兴起，那就是假基督的角。它要在它之前拔除那三个；也就是说，他要推翻\Place{埃及}、\Place{利比亚}和\Place{埃塞俄比亚}的三王，目的是为自己获取普世主权。在征服其余七角之后，他最终将开始，被一种奇怪而邪恶的灵所鼓动，向圣人挑起战争，到处迫害所有人，意图被所有人荣耀，并被崇拜为神。\par

22. \Quote{直到万古常存者来到。} 那就是，当最终审判者的审判者和万王之王的从天降临时，他将推翻仇敌的全部统治和权能，并用永恒的惩罚之火烧灭一切。但给他的仆人们、先知们、殉道者们和所有敬畏他的人，他要赐予一个永远的国；也就是说，他们将拥有无尽的善的享受。\par

25. \Quote{直到一载、二载、半载。} 这表示三年半。\par

\BibleRef{达 9:21} \Quote{看，加俾额尔那人……飞着。} 你看先知如何将天使的速度比作有翼的鸟，因为这种灵体在执行命令时飞翔得如此迅速，其运动轻盈而快速。\par

\BibleRef{达 10:6} \Quote{祂话语的声音。} 因为所有如今信靠他的人，都声明基督的话，仿佛我们通过他的口说出他所命令的事。\par

7. \Quote{我看见，} 等等。因为他是向敬畏他的圣人，唯独向他们启示自己。如果有人在教会中似乎现在活着，却没有敬畏天主的心，他与圣人的交往对他毫无益处。\par

12. \Quote{你的话已蒙垂听。} 看哪，一个义人的虔敬何等重要，唯独向他，如同向一个配得的人，世上尚未显现的事被启示出来。\par

13. \Quote{看，弥额尔。} \Person{弥额尔}是谁？不就是派给那民族的天使吗？正如（天主）对\Person{梅瑟}说：\Quote{我不同你上去，因为百姓是顽固的；但我的天使要同你上去。}\par

16. \Quote{我的心肠翻转}（《英王钦定本》作\InnerQuote{我的忧愁临到我}）。因为理当在主显现时，上方的要转到下方，以便下方的也能升到上方。——他说，我需要时间恢复自己，才能承受那些话并回答所说的。——但他继续说，当我在这个状态时，我出乎意料地被加强了。因为一位看不见的触摸了我，立刻我的软弱被除去，我恢复到了从前的力量。因为当我们生命的一切力量和荣耀离开我们时，我们就靠着基督被加强，他伸出他的手，从死者中复活活着的人，就像从阴府本身，升至生命的复活。\par

18. \Quote{他就加强了我。} 因为每当\OriginalTerm{圣言}{Word}使我们对未来充满希望时，我们也能轻易听见他的声音。\par

20. \Quote{与波斯国的君长争斗。} 因为从你在主你的天主面前谦卑自己的那一天起，你的祈祷就被垂听，我被派遣\Quote{与波斯国的君长争斗}。因为有一个计划不让他们的人民离开。因此，为使你的祈祷迅速蒙应允，\Quote{我起来对抗他。}\par

\BibleRef{达 12:1} \Quote{将有一个灾难时期。} 因为那时将有大的灾难，如同世界奠基以来从未有过的，那时一些人以这种方式，另一些人以那种方式，被派遣到各城各国去毁灭信众；圣人将从西方旅行到东方，从东方被驱赶到南方，受迫害；其他人将藏在山中和洞穴里；那可憎之物要处处攻击他们，通过他的法令在海中和陆地上剪除他们，并设法用一切手段将他们从世上消灭；他们将不能再出卖自己的财产，也不能从陌生人那里购买，除非有人保留并随身带着那兽的名号，或在额上带着它的印记。因为那时他们将被从各处驱逐，从自己的家中被拖出，被拉进监狱，受各种刑罚，并从全世界被赶出。\par

2. \Quote{这些人要醒起而得永生。} 那就是，那些相信真生命的人，他们的名字写在\OriginalTerm{生命册}{book of life}上。\Quote{这些要蒙羞。} 那就是，那些依附假基督的人，他们与他一起被投入永远的惩罚中。\par

3. \Quote{智者要发光。} 主在福音中也说了同样的话：\Quote{那时义人要在他们父的国里发光如同太阳。}\par

7. \Quote{一载、二载、半载。} 借此他指明了假基督的三年半。因为他用\Quote{一段时间}表示一年；\Quote{多段时间}表示两年；\Quote{半段时间}表示半年。这就是但以理预言的那\Quote{一千二百九十日}。\par

9. \BibleQuote{言语已经封闭和密封。}因为人无法说明天主为圣人们预备了什么；因为眼未曾看见，耳未曾听见，人心也未曾想到（领会）这些事，连圣人们到那时也将切望观看；所以祂对他说：\BibleQuote{因为言语的封闭，直到末期；直到许多人被拣选并受火试炼。}谁是那被拣选的呢？岂非那些相信真理之言，因而得以洁白，脱去罪的污秽，穿上那属天的、纯洁的、荣耀的圣神，以便在新郎来时，他们可以与他一同进去？\par

11. \Quote{那行毁坏的可憎之物将被放置（设立）。}因此，达尼尔论及两种可憎之物：一种是毁坏的可憎之物，安提约古在定的时候设立了它，它与那荒凉的可憎之物相关；另一种是普世性的，即当假基督来临时。因为，如达尼尔所说，他也将被设立，为毁灭许多人。\par

% structure: p0259 source=h2 target=section reason=relative-heading-rank; base=h2

\section{论达尼尔的其他片段}

因为当现今持有王权的铁腿让位于脚和脚趾时，按那可怕兽的象征，正如先前所预示的，那时从天上会来一块击打像并打碎它的石头；它将颠覆所有的国度，并将国度交给至高者的圣徒。这就是那块变成一座大山、充满大地的石头，经上论及它说：\BibleQuote{我在夜间的异象中观看，见有一位像人子的，驾着天云而来，来到亘古常在者面前。祂得了权柄、荣耀、国度，使各方、各国、各族的人都事奉祂；祂的权柄是永远的，不能废去，祂的国度也不被毁灭。}\par

% structure: p0261 source=h3 target=subsection reason=relative-heading-rank; base=h2

\subsection{论三青年之歌}

\BibleQuote{阿纳尼亚、阿匝黎雅、米沙耳，请赞美主；主的宗徒、先知和殉道者，请赞美主：赞美祂，尊崇祂超越一切，直到永远。}\par

我们实在惊叹于火窑中三位青年的言语，他们列举了所有受造物，好使它们中没有一个可被视为独立自存的；而是总括并称呼它们全部，无论在天上、地上、地下的，都表明它们都是天主的仆人，祂藉着圣言创造了万有，以免有人夸耀任何受造物是无生无始的。\par

% structure: p0264 source=h3 target=subsection reason=relative-heading-rank; base=h2

\subsection{论苏撒纳}

这里叙述的事发生在较晚的时候，尽管它被放在第一卷（书卷开头）之前。因为作者们常以颠倒顺序在他们的著作中叙述许多事情。我们在先知书中也发现，有些异象记录在先，却应验在后；反之，另一些记录在后，却应验在先。这是出于圣神的安排，使魔鬼不能理解先知以比喻所说的事，免得他再次设下罗网毁灭人。\par

1节。\Quote{名叫约雅金。}这约雅金在巴比伦作客，娶了苏撒纳为妻。她是司铎赫耳基亚的女儿，赫耳基亚在约史雅王命他洁净至圣所时，在主的殿中找到了法律书。他的兄弟是耶肋米亚先知，他在百姓被掳到巴比伦后，与留下的遗民一同被带往埃及，住在塔弗乃；在那里宣讲预言时，被百姓用石头砸死。\par

\Quote{极其美貌，且敬畏主，}等等。因为从结出的果实，树也易于被认识。因为虔诚而热心于法律的人，将带出配得天主的儿女；就如那位成为基督的先知和见证人，以及那位在巴比伦被发现贞洁而忠信的女子，她的尊荣与贞洁成了有福的达尼尔彰显为先知的机缘。\par

4节。\Quote{约雅金非常富有，}等等。因此我们须寻求解释。因为那些被掳并臣服于巴比伦人的人，怎能像在自己家中一样在同一地方聚集？对此我们应观察到，拿步高在将他们迁离后善待他们，并允许他们聚集，且按法律行一切事。\par

7节。\Quote{到了正午，苏撒纳进入（她丈夫的花园）。}苏撒纳预表了教会；她的丈夫约雅金预表了基督；花园预表了诸圣的蒙召，他们像结果实的树被栽种在教会中。巴比伦预表了世界；两位长老被呈现为那图谋反对教会的两个民族的象征——即受割礼的与外邦人。因为经上说\Quote{被立为民的首领和审判者}，意思是在这世上他们行使权柄并统治，不义地审判义人。\par

8节。\Quote{两位长老看见了她。}犹太人的首领如今想要从书中删除这些事，并断言这些事未在巴比伦发生，因为他们对那时长老所做之事感到羞愧。\par

9节。\Quote{他们扭曲了自己的心。}因为那些曾是教会仇敌和败坏者的人，怎能以纯洁的心公义审判或仰望上天呢？他们已成了这世界之君的奴隶。\par

10节。\Quote{两人都被（她的爱）所伤。}这话要按真实意义理解；因为这两个民族常被在他们内运行的撒殚所伤（激动），力图对教会兴起迫害和磨难，并寻求如何败坏她，尽管他们彼此并不一致。\par

12节。\Quote{他们殷勤地窥探。}这一点也值得注意。因为直到如今，外邦人和受割礼的犹太人都留心并忙于教会的事务，渴望针对我们收买假证人，如宗徒所说：\BibleQuote{因为有偷着引进来的假弟兄，他们私下窥探我们在基督耶稣内所有的自由。}\par

将心思放在女人身上是一种罪过。\par

14. \Quote{他们出去以后，彼此分离。}至于他们在午餐时分彼此分离，这表明在属世的肉食上，\OriginalTerm{犹太人}{Jews}和\OriginalTerm{外邦人}{Gentiles}并不一致；但在他们的见解和一切世俗事务上，他们却同心合意，能够彼此相遇。\par

14. \Quote{他们彼此询问，就承认了自己的情欲。}因此，他们彼此揭露，预表了那时刻：人将凭自己的心思受到考验，并要为所行的一切\OriginalTerm{罪}{sin}向\OriginalTerm{天主}{God}交账，正如\Person{撒罗满}{Solomon}所说：\Quote{审查必毁灭不敬虔的人。}因为这些人是被审查定罪的。\par

15. \Quote{他们窥伺适当的时候。}这适当的时候若不是\OriginalTerm{逾越节}{Passover}，又是什么呢？那时，在花园中为那些焚烧的人预备了洗濯盆，\Person{苏撒纳}{Susannah}洗浴自己，并作为纯洁的新娘被献给\OriginalTerm{天主}{God}。\par

\Quote{只有两个使女。}因为当\OriginalTerm{教会}{Church}愿意按照惯例取用洗濯盆时，她必须有两个使女陪伴。因为教会是藉着对\OriginalTerm{基督}{Christ}的\OriginalTerm{信德}{faith}和对\OriginalTerm{天主}{God}的\OriginalTerm{爱德}{love}宣认并领受洗濯盆的。\par

18. \Quote{她对使女说：给我拿油来。}因为\OriginalTerm{信德}{faith}和\OriginalTerm{爱德}{love}为那些被洗的人预备\OriginalTerm{油}{oil}和\OriginalTerm{香膏}{ointments}。但这些香膏是什么，岂不就是神圣\OriginalTerm{圣言}{Word}的\OriginalTerm{诫命}{commandments}？油是什么，岂不就是\OriginalTerm{圣神}{Holy Spirit}的能力，信徒在洗濯之后，如同用香膏被傅抹？这一切都在有福的\Person{苏撒纳}{Susannah}身上预先描绘，为了我们的缘故，好使我们这些现今信\OriginalTerm{天主}{God}的人，不以现今在\OriginalTerm{教会}{Church}中所行的事为奇怪，而相信这一切都是由古圣祖以预象所表明的，正如\OriginalTerm{宗徒}{apostle}也说：\BibleQuote{这些事发生在他们身上，作为预表；并且写下来是为了教训我们这处于世界末期的人。}\par

18. \Quote{他们从暗门出去；}如此预先表明，凡愿意在花园中分享\OriginalTerm{水}{water}的人，必须弃绝宽门，而从窄门进入。\par

\Quote{他们没有看见长老。}因为正如古时\OriginalTerm{魔鬼}{devil}藏在花园中的蛇里，现在也同样藏在长老们里面。他用他自己的\OriginalTerm{情欲}{lust}点燃他们，好使他再次败坏\Person{厄娃}{Eve}。\par

20. \Quote{看，花园的门都关了。}邪恶的统治者，充满\OriginalTerm{魔鬼}{devil}的作为，\Person{梅瑟}{Moses}将这些事传给你们吗？你们自己诵读法律，就这样教导别人吗？你们说：\BibleQuote{不可杀人}，却杀人吗？你们说：\BibleQuote{不可贪恋}，却贪恋邻人的妻子吗？\par

\Quote{我们爱上了你。}你们这些无法无天的人，为什么用欺骗的话试图赢得一个贞洁无邪的\OriginalTerm{灵魂}{soul}，以满足你们自己的\OriginalTerm{情欲}{lust}？\par

21. \Quote{若你们不肯，我们就要作证控告你们。}你们开始时的这种邪恶胆量，来自起初就潜伏在你们里面的诡诈。而且确实有一个青年人与她同在，那是属于你们的一位；一个从天上来的人，不是要与她交合，而是要为她的\OriginalTerm{真理}{truth}作证。\par

22. \Quote{\Person{苏撒纳}{Susannah}叹息了。}有福的\Person{苏撒纳}{Susannah}听见这些话，心中烦乱，在嘴上设了守卫，不愿被邪恶的长老玷污。现在我们也能领悟发生在\Person{苏撒纳}{Susannah}身上之事的真实含义。因为你们可以发现这事也在\OriginalTerm{教会}{Church}的现今状况中应验了。因为当两个民族合谋毁灭任何圣人时，他们窥伺适当的时候，进入\OriginalTerm{天主}{God}的殿，此时所有人都在祈祷赞美\OriginalTerm{天主}{God}，他们抓住其中一些人，带走并扣留他们，说：\Quote{来，与我们同流合污，敬拜我们的神；否则，我们就要作证控告你们。}当他们拒绝时，就把他们拖到法庭，控告他们违反了\Person{凯撒}{Caesar}的法令，并判处他们死刑。\par

\Quote{我四面受困。}看哪，一个贞洁且为\OriginalTerm{天主}{God}所爱的女人的话：\Quote{我四面受困。}因为\OriginalTerm{教会}{Church}受折磨和困窘，不仅来自\OriginalTerm{犹太人}{Jews}，也来自\OriginalTerm{外邦人}{Gentiles}，还来自那些称为基督徒、实际上并非如此的人。因为他们看见她贞洁幸福的生活，就竭力要毁灭她。\par

\Quote{因为我若做这事，就是死亡。}因为不顺从\OriginalTerm{天主}{God}而顺从人，必导致永死和惩罚。\par

\Quote{我若不做这事，也不能逃脱你们的手。}这话确实说得真实。因为那些为\OriginalTerm{天主}{God}的名受审判的人，若顺从人的命令，就在\OriginalTerm{天主}{God}面前死了，却要在世界活着。但他们若拒绝顺从人的命令，就逃脱不了审判官的手，反而被他们定罪。\par

\Quote{我宁可不做。}因为死于恶人之手而与\OriginalTerm{天主}{God}同活，远胜于顺从而从他们手中得释放，却落在\OriginalTerm{天主}{God}的手中。\par

24. \Quote{\Person{苏撒纳}{Susannah}大声呼喊。}\Person{苏撒纳}{Susannah}向谁呼喊，岂不是向\OriginalTerm{天主}{God}吗？正如\Person{依撒意亚}{Isaiah}所说：\BibleQuote{那时你呼求，\OriginalTerm{主}{Lord}必应允你；你还在说话，祂必说：我在这里。}\par

\Quote{两个长老就高声反对她。}因为恶人从不停止反对我们，并说：\Quote{从地上除掉这样的人，因为他们不配活着。}在福音的意义上，\Person{苏撒纳}{Susannah}藐视那些杀身体的人，为要救她的\OriginalTerm{灵魂}{soul}脱离死亡。\OriginalTerm{罪}{sin}是灵魂的死亡，尤其是\OriginalTerm{奸淫}{adultery}。因为当与\OriginalTerm{基督}{Christ}结合的\OriginalTerm{灵魂}{soul}离弃了她的\OriginalTerm{信德}{faith}，它就被交付于永死，即\OriginalTerm{永罚}{eternal punishment}。为证实这事，在肉体的婚约违犯和破坏的案例中，法律也规定了死刑。\par

\Quote{于一人跑去开了门；}指向那宽大之路，跟随这类人的人便由此灭亡。\par

31. \Person{苏撒纳}{Susannah}是个极其柔弱的女子。不是像\Person{依则贝耳}{Jezebel}那样有妖艳的妆饰，或眼涂各色颜料；而是有\OriginalTerm{信德}{faith}、\OriginalTerm{贞洁}{chastity}和\OriginalTerm{圣洁}{sanctity}的妆饰。\par

\Quote{他们把按手在她头上}，好至少藉触摸她来满足他们的情欲。\par

\Quote{她哭泣着。}因为她的眼泪从天上吸引了圣言的注意，祂将要以眼泪复活死去的拉匝禄。\par

\Quote{会众就相信了他们。}因此，我们理当在每项职责上坚定不移，不理会谎言，也不对统治者的身份曲意服从，知道我们须向天主交账；但若我们追随真理，持守信德的准确准则，我们必蒙天主悦纳。\par

\Quote{上主垂听了她的声音。}因为凡以纯洁之心呼求祂的，天主必垂听；但那以诡诈和伪善（呼求祂）的，天主必掩面不顾。\par

\Quote{你这在邪恶中衰老的。}既然我们在开头的引言中已经说明，那两位长老是预表两个民族，即受割损的和外邦人的民族，他们常是教会的仇敌；现在让我们注意达尼尔的话，并明白圣经绝没有欺骗我们。因为祂对第一位长老说话，谴责他，视他为一个受过法律教导的人；而对另一位长老，祂视他为外邦人，称他为\Quote{客纳罕的种子}，尽管他当时属于受割损的人。\par

\Quote{因为即使现在，天主的使者。}祂也显示，当苏撒纳向天主祈祷并蒙垂听时，天使就在那时被派去帮助她，正如在多俾亚和撒辣的事例中一样。因为他们祈祷时，他俩的恳求在同一天同一时辰蒙了垂听，辣法耳天使被派去医治他们二人。\par

\Quote{他们就起来攻击那两个长老；}为要应验这句话：\Quote{谁给邻舍掘坑，必掉在其中。}\par

因此，亲爱的诸位，我们理当留心这一切，恐怕有人陷入任何过犯，冒失去灵魂的危险，因为我们知道天主是众人的审判者；圣言本身就是那眼睛，世上所行的一切都不能逃过它。因此，让我们常以警醒的心和纯洁的生活，效法苏撒纳。\par

% structure: p0302 source=h2 target=section reason=relative-heading-rank; base=h2

\section{论玛窦福音}

\BibleRef{玛 6:11}\par

为此，我们被命令求取足以维持身体实质的东西：不是奢侈，而是食物，它恢复身体所失去的，并防止因饥饿而死亡；不是煽动并驱向享乐的宴席，也不是那些使肉体悖逆灵魂的东西；而是面包，而且不是为多年，而是为今天够用的。\par

% structure: p0305 source=h2 target=section reason=relative-heading-rank; base=h2

\section{论路加福音}

\BibleRef{路 2:7} 若你愿意，我们说圣言是天主首生者，祂从天上降到荣福玛利亚那里，并在她腹内成为首生的人，为要使天主的首生者与首生的人结合而显现。\par

\BibleRef{路 2:22} 当他们带祂上耶路撒冷，将祂献给主时，他们献上了取洁的祭品。因为若为祂献上按法律所需的取洁礼物，那么祂实在成了法律之下的。但圣言并不像那些谄媚者所幻想的那样服从法律，因为祂本身就是法律；天主也不需要取洁的祭品，因为祂在片刻之间洁净并圣化一切。然而，祂取了从童贞女所接受的人性，成为法律之下的，并像首生者一样受取洁礼，并非因为祂需要这种仪式而履行其礼仪，而是为了救赎那些因咒诅的审判而被卖于法律束缚之下的人。\par

\BibleRef{路 23} 为此，阴府的门卫看见祂时战栗；铜门和铁闩也被折断。因为，看，独生子进入了，灵魂中的灵魂，天主圣言带着（人的）灵魂。因为祂的身体躺在坟墓中，并没有失去天主性；正如祂在阴府时，在本质存在上与祂的父同在，同样祂也在身体内并在阴府内。因为圣子不受空间限制，如同圣父一样；祂包容万有在祂自己内。但祂凭自己的意愿居住在由灵魂赋予生命的身体里，为的是祂能以祂的灵魂进入阴府，而不是以祂纯粹的天主性。\par

% structure: p0309 source=h2 target=section reason=relative-heading-rank; base=h2

\section{关于梅瑟五书存疑残篇}

% structure: p0310 source=h3 target=subsection reason=relative-heading-rank; base=h2

\subsection{序言}

因父、及子、及圣神之名，独一的天主。这是卓越的法律的抄本。但在开始缮写法律书的抄本之前，兄弟啊，值得先就它的卓越和其安排的高贵来教导你。它的首要卓越之处是：天主藉我们最蒙福的领袖、先知之首和第一位宗徒，即被派往以色列子民的那位——阿默兰的儿子、刻哈特的孙子、肋未的后裔梅瑟——的手交付了它。他装饰有一切智慧，赋有最好的天才。他尊严显赫，性格正直非凡，推理能力出众，与天主交谈。天主拣选他作为贵重的器皿。至高天主藉祂的领袖和先知将它传给了我们，并信托给了我们（愿祂的名受赞美），使用的是塔尔古木的叙利亚语，七十贤士将它译为希伯来语，即民族的言语和普通百姓的惯用语。梅瑟遂从永恒的主接受了它，是第一个受托付并遵守其规条和法令的人。然后他将它教导给以色列子民，他们也接受了它。他向他们解释了其中深奥的奥秘和隐晦之处。他也向他们阐释了那些较不易理解的事，正如天主所允许的，并隐藏了那些天主禁止他（启示）的法律秘密。在他们中间，没有一人像他那样在祂的审判和法令上更熟练，并更清晰地传述祂教义的奥秘，直到天主在旷野中使他经过整整四十年的成熟后，将他接到自己那里。\par

以下是继梅瑟先知之后，连续传承法律直到默西亚来临的教师们之名：\par

那么，我兄弟，愿天主祝福他，应当知道：天主将最卓越的法律交托给先知\Person{梅瑟}，\Person{阿默兰}之子。\par

\Person{厄斯德拉}将它交托给大司祭\Person{沙迈}，\Person{雅杜亚}交托给\Person{撒米安}，\Person{撒米安}交托给\Person{安提哥诺}。\par

\Person{西默盎}交托给\Person{耶胡达}。\par

\Person{耶胡达}交托给司祭\Person{匝加利亚}。\par

\Person{若瑟}交托给\Person{哈南}和\Person{盖法}。此外，司祭职、君王职和先知职从他们那里被夺去。\par

这些人在默西亚来临之际是教师；他们二人也是以色列子民的司祭。因此，受托保管这最卓越法律的可敬和尊贵司祭总数是五十六位，\Person{哈南}（即\Person{亚纳斯}）和\Person{盖法}除外。\par

这些就是在末期将它交托给以色列子民团体的人；在他们之后没有再兴起任何司祭。\par

这就是关于最卓越法律所发生之事的记述。\par

\Person{阿尔米乌斯}，\WorkTitle{时代}一书的作者曾说：在\Person{托勒密}王在位第十九年，他命令召集以色列子民的长老，好让他们将法律的一个抄本交给他，并且每人能随时解释其含义。\par

于是长老们来了，带来最卓越的法律。然后他命令他们每个人都为他解释法律书。\par

但他不同意长老们所作的解释。他命令将长老们投入监狱和锁链中。他抓住法律书，把它扔进深坑，并在其上投放火和热灰达七天之久。然后他命令他们将城中的污秽扔进那有法律书的坑里。坑被填满了，直到顶端。\par

法律在那坑的污秽下存留了七十年，却没有朽坏，甚至没有一片叶子损坏。\par

在\Person{阿皮亚努图斯}王在位第二十一年，他们将法律书从坑中取出，没有一片叶子损坏。\par

在基督升天之后，罗马王\Person{阿斯帕西亚努斯}之子\Person{提图斯}王来到\Place{耶路撒冷}，围攻并攻取了它。他摧毁了以色列子民所建造的第二圣殿的建筑。\Person{提图斯}王摧毁了圣所之殿，杀死了其中所有犹太人，并用他们的血建造了\Place{熙雍}（\Emph{sic}）。在那次掳掠之后，犹太人被分散到各地为奴。他们再也没有在\Place{耶路撒冷}城聚集，也没有任何返回的希望。\par

因此，\Place{耶路撒冷}被荒废之后，\Person{舍玛亚}和\Person{安塔利亚}（\Person{阿布塔利翁}）将法律交托出去——他们是\Place{巴阿尔巴克}的王，这城原是\Person{达味}王之子\Person{索里曼}古时所建，并在\Person{默纳舍}王的日子重新修复，\Person{默纳舍}曾将\Person{依撒意亚}先知锯成两半。\par

厄东子孙的\Person{哈德良}王围攻\Place{巴阿尔巴克}，攻取了它，杀了其中所有犹太人，并将所有\Person{达味}家族的人掳为奴隶。犹太人被分散在全地，正如至高天主所预言的：\Quote{我要将你们分散在外邦人中，将你们分散在列国中。}\par

这些就是关于那最卓越之书的历史所传达到我们的事。前言结束。\par

% structure: p0330 source=h3 target=subsection reason=relative-heading-rank; base=h2

\subsection{法律}

因永恒、永存、至能、仁慈、怜悯的天主之名。\par

赖天主的助佑，我们开始描述法律书及其诠释，正如神圣、博学、最卓越的教父们所诠释的。\par

因此，以下就是第一卷书的诠释，它确实是受造物（和）创造之书。\par

% structure: p0334 source=h3 target=subsection reason=relative-heading-rank; base=h2

\subsection{第一节}

论天地之创造。\Quote{在起初天主创造了}等等。\par

对天主所说之事的阐述。\par

确实，蒙福的先知，伟大的\Person{梅瑟}，写了这本书，并指定和标记其标题为\Emph{存在之书}，即\Quote{论受造物}等等。\par

% structure: p0338 source=h3 target=subsection reason=relative-heading-rank; base=h2

\subsection{第二、三节}

第二、三节。主说：\Quote{我要使洪水的水临到地上，毁灭所有血肉}等等。\par

\Person{希波利图斯}，塔尔古玛释义者，说：诺厄儿子们的妻子名字如下：\Person{闪}的妻子名叫纳哈拉特·玛赫努克；\Person{含}的妻子名叫则德卡特·纳布；\Person{雅斐特}的妻子名叫阿拉特卡。此外，这些是他们在叙利亚塔尔古玛中的名字。\Person{闪}的妻子名叫纳哈拉特·玛赫努克；\Person{含}的妻子名叫则德卡特·纳布；\Person{雅斐特}的妻子名叫阿拉特卡。\par

因此天主向\Person{诺厄}提示，通知他洪水的来临，以及败坏（邪恶）之人的毁灭。\par

至高天主命令他，他和他的儿子们以及儿子们的妻子，从圣山下来，建造一艘三层的大船。下层是留给凶猛、野生、危险的野兽。它们之间有木桩或木梁隔开，以防止它们互相接触。中层是留给鸟类的，以及它们不同的物种。上层则是留给\Person{诺厄}本人和他的儿子们——以及他自己的妻子和儿子们的妻子。\par

\Person{诺厄}还在大船的东侧造了一扇门。他还建造了水箱和储藏食物的仓库。\par

因此，当\Person{诺厄}造完大船后，他与儿子们——\Person{闪}、\Person{含}和\Person{雅斐特}——进入了寄存的洞穴。\par

在他们初次走近时，他们确实欣喜地找到了先祖们的遗体：\Person{亚当}、\Person{舍特}、\Person{厄诺士}、\Person{刻南}、\Person{玛哈拉耳}、\Person{雅勒得}、\Person{玛突舍拉}和\Person{拉默客}。这八具遗体在寄存之处，即\Person{亚当}、\Person{舍特}、\Person{厄诺士}、\Person{刻南}、\Person{玛哈拉耳}、\Person{雅勒得}、\Person{玛突舍拉}和\Person{拉默客}的遗体。\par

此外，\Person{诺厄}取走了\Person{亚当}的遗体。他的儿子们随身携带了祭品：\Person{闪}带了黄金，\Person{含}带了没药，\Person{雅斐特}带了乳香。然后，他们离开寄存的洞穴，将祭品和\Person{亚当}的遗体转移到圣山。\par

当他们坐在\Person{亚当}的遗体旁，面对乐园时，他们开始哀悼哭泣，为失去乐园而悲伤。\par

然后，他们从圣山下来，举目望向乐园，重新开始哭泣哀号，并用这些话作了永远的告别：\Quote{永别了！愿平安归于你，天主之乐园！永别了，宗教与纯洁的住所！永别了，快乐与愉悦的宝座！}\par

于是他们拥抱圣山上的石头和树木，哭泣说：再会了，善者的居所！再会了，圣善身体的住所！\par

三天后，\Person{诺厄}偕同儿子们和儿媳们从圣山下来，到圣山的山脚，到方舟所在之处。因为方舟是在圣山的突出边缘之下。\par

\Person{诺厄}进入方舟，将\Person{亚当}的遗体与供物放在方舟中央，安放在他为接纳遗体所预备的木架上。\par

天主吩咐\Person{诺厄}说：你要为自己用黄杨木（或柏木）做成响板。\par

也要用同一木料制作其锤（钟）。响板的长度应为三整肘，宽度应为肘半。\par

天主命令他每日击打响板三次：第一次在黎明时分，第二次在正午，第三次在日落时分。\par

每当\Person{诺厄}击打响板，\Person{加音}的儿子们和\Person{瓦希姆}的儿子们立刻跑到他那里；他便警告并惊吓他们，告知洪水即将来临，毁灭已迫在眉睫。\par

这样，天主向他们显示怜悯，使他们悔改并清醒。但\Person{加音}的儿子们不听从\Person{诺厄}向他们宣告的话。\Person{诺厄}将各类飞鸟、禽兽，驯良的和野生的，都一对一对地集合起来。\par

% structure: p0357 source=h3 target=subsection reason=relative-heading-rank; base=h2

\subsection{第四节}

论\BibleRef{创 7:6}。\par

\OriginalTerm{塔古木}{Targum}的叙利亚注释者\Person{希坡律陀}说：我们在一份古老的希伯来抄本中发现，天主命令\Person{诺厄}将野兽依次排列在下层，并用木桩将雄性分开雌性。\par

同样，他也这样对待所有牲畜，以及中层的飞鸟。天主命令将雄性分开雌性，是为了庄重和洁净，以免它们彼此混杂。\par

此外，天主对\Person{梅瑟}说：为你和你的儿女预备食物。这些食物要用小麦磨碎、捣烂、用水调和并晒干。\Person{诺厄}当即吩咐他的妻子和儿媳们专心揉面并放进炉中。她们便揉面，预备仅够他们所需的数量，以致几乎没有剩余。\par

天主吩咐\Person{诺厄}说：谁先向你报告洪水来临，你立刻将他消灭。当时\Person{含}的妻子正站着，要把一大块面包放进炉中。忽然，照主的话，水从炉中涌出，水流穿透并毁坏了面包。因此\Person{含}的妻子对\Person{诺厄}说：先生，天主的话实现了：\Quote{天主预言的已经成就；}所以，请执行主所吩咐的。\Person{诺厄}听见\Person{含}妻子的话，便对她说：难道洪水已经来了吗？\Person{含}的妻子对他说：你已经说了。天主却立刻吩咐\Person{诺厄}说：不要消灭\Person{含}的妻子；因为从你口中说出了毁灭的开始——\Quote{你先说洪水来了。}因\Person{诺厄}的声音洪水来到，水突然毁灭了那面包。天上的水闸被打开，雨水倾泻在地上。那最初说过\BibleQuote{让水聚在一处，使旱地露出来}的同一声音，准许众水的泉源和海涛自行涌出，并带来了大水。\par

试想天主关于世界所说的话：让它的所有高处降低，它们就降低了；让它的低处从深处升起。\par

大地被剥去并空无一物，如同起初一样。\par

雨水从上倾注，地底涌出；大地的框架被毁，原始秩序被打破。世界变得如同起初被水淹没时那样荒凉。其中的任何存在物都未保持完整。\par

它原先的结构遭到破坏，大地被骤然涌来的洪水和其漫溢的规模、众多的暴雨以及从深处的喷发所损毁，因为水不断涌出。最终，它变得如同从前一样。\par

% structure: p0367 source=h3 target=subsection reason=relative-heading-rank; base=h2

\subsection{第五节}

论\BibleRef{创 8:1}。\par

\OriginalTerm{塔古木}{Targum}的注释者\Person{希坡律陀}和我的导师\Person{雅各伯·罗哈维安西斯}说：在希伯来历二月即亚珥月第二十七日，方舟从圣山脚下升起；水已托起它，载着它在世界的四个基本方向漂流。方舟从圣山向东离开，然后折向西，又转向南，最后向东航行，于第十个月的第一日靠近\Place{卡尔杜山}。那便是坎努月第二个月。\par

\Person{诺厄}在第二年的亚珥月第二十七日出了方舟：因为方舟航行了整整五个月，在水上移动往来，并在五十一天内靠近陆地。此后它不再漂浮，只是依次朝向地的四个基本方向移动，最终又稳定地朝向东方。此外，我们说，那是十字架的标记。方舟是预表将要来临的基督。因为那方舟拯救了\Person{诺厄}和他的儿子们，也拯救了牲畜、走兽和飞鸟。同样，基督在十字架上受苦时，拯救了我们脱离控诉和罪恶，并用祂最纯洁的鲜血洗净了我们。\par

正如方舟回到东方并靠近\Place{卡尔杜山}，同样，基督当祂预定的工作完成时，回到天上圣父的怀中，坐在圣父右边荣耀的宝座上。\par

至于卡尔杜山，位于东方，在拉班子孙之地；东方人称它为哥达什山，阿拉伯人和波斯人称它为亚拉腊山。\par

另有一座名为卡尔杜的城，该山便以它命名；山势极高又无法攀登，因常有狂风暴雨，从来无人能到达山顶。若有人试图攀登，便有魔鬼冲来，将他从山脊摔下平原，令他死亡。此外，无人知晓山顶有何物，唯有某些方舟木料的遗物仍散落在山顶表面。\par

% structure: p0374 source=h3 target=subsection reason=relative-heading-rank; base=h2

\subsection{第十节}

第十节。论\BibleRef{申 33:2}。\par

希波利特，\OriginalTerm{塔古姆}{Targum}的解释者，曾说：梅瑟在完成这预言之后，也按各支派祝福了所有以色列子孙，并为他们祈祷。然后天主吩咐梅瑟说：\Quote{上乃波山去，那山在摩阿布地，面对耶里哥，希伯来人称之为希伯来人之山。}\par

祂对他说：\Quote{你观看客纳罕地，就是我要赐给以色列子孙为产业的地方。然而你决不可进去，所以你从远处好好观看。}梅瑟观看时，看见那地——一片青绿、富饶丰盛、树木茂密；梅瑟大为感动，哭泣起来。\par

梅瑟从乃波山下来，召来农的儿子若苏厄，在以色列子孙面前对他说：\Quote{你要刚强勇敢，因为你将带领以色列子孙进入天主应许赐予他们祖先为产业的地方。所以不要惧怕百姓，也不要畏惧异民，因为天主必与你同在。}\par

梅瑟写了那部\OriginalTerm{森纳}{Senna}（希伯来文\par

＝\Quote{次级法律}，或\Quote{申命纪}），交给肋未子孙的司祭，吩咐他们说：\Quote{七年内隐藏这部森纳，在整整七年内不可出示。（\InnerQuote{然后}）在帐棚节上，肋未子孙的司祭要在以色列子孙面前诵读这法律，使全体人民，男女都遵守天主的话：吩咐他们遵守那法律中天主的话。凡违犯其中一条诫命的，当受诅咒。}\par

于是，梅瑟写完法律后，将它交给农的儿子若苏厄，并命令他将法律交给肋未子孙的司祭。梅瑟又叮嘱他们，要将法律书再放入主的约柜内，使它永远留在那里作为证物。\par

梅瑟完成一切吩咐后，天主命他上乃波山，山在耶里哥对面。主将整个应许之地从四方指给他看：从旷野到海，从海到海。主对他说：\Quote{你亲眼看见了，但你决不可进去。}于是，天主的仆人梅瑟，按天主的命令死在那里。天使将他埋葬在乃波山上，面对贝特培敖尔。直到今日，无人知晓他的坟墓，因为天主隐藏了他的墓地。\par

梅瑟活了一百二十岁，他的眼睛没有昏花，脸上的皮肤也没有皱纹。\par

梅瑟在某日死去，那日是第二个月第七日，即依雅尔月，在一天的第三时辰。\par

以色列子孙在摩阿布平原为他哀哭了三日。\par

农的儿子若苏厄充满了智慧之神，因为梅瑟曾按手在他头上。全体以色列子孙都听从了他。天主在某日——即尼散月第七日——吩咐了农的儿子若苏厄。\par

农的儿子若苏厄活了一百一十岁，在第四日死去，那是厄路耳月的第一日。他们将他葬在厄弗辣因山上的塔姆纳特塞辣城。\par

愿天主因\Emph{著作的完成}而受赞美。\par

% structure: p0389 source=h2 target=section reason=relative-heading-rank; base=h2

\section{论圣咏}

论圣咏。一、罗马主教希波利特所撰\WorkTitle{圣咏集释义}的绪论。\par

1. \WorkTitle{圣咏集}在梅瑟所赐的法律之后，包含了新的教义；因此它是梅瑟圣经之后的第二部教义书。梅瑟和若苏厄死后，在民长之后，达味兴起了，他被视为配得称为救主之父的人；他是第一位给希伯来人带来新式圣咏的人，借此废除了梅瑟关于祭献的规条，并引入了以赞美诗和欢呼敬拜天主的新方式；此外，他在整个职务中也教导了许多超越梅瑟法律的事。这便是这书的神圣性和益处。关于它的题名，须说明如下：由于许多在基督内相信的弟兄认为这书是达味的，并称之为\WorkTitle{达味圣咏}，我们须陈述所流传给我们的内容。希伯来人称这书为\OriginalTerm{塞弗拉·特林}{Sephra Thelim}，而在\WorkTitle{宗徒大事录}中它被称为\WorkTitle{圣咏集}（经文这样说：\Quote{正如\WorkTitle{圣咏集}上所记载的}），但书的题名中并未出现作者的名称。原因是，那里所写的不是一个人的话，而是数人合在一起的话；据传统，厄斯德拉在充军之后，将数人的圣咏，或更确切地说，将他们的话汇编成一卷，因为它们并非全是圣咏。因此有些冠以达味之名，有些冠以撒罗满之名，有些冠以阿撒夫之名。另有一些属于伊狄通（耶杜通）；此外还有一些属于科辣黑之子，甚至属于梅瑟。既然是这样多人的话汇集在一起，任何明白事理的人都不会说它们仅出自达味一人。\par

2. 至於那些沒有題署的，我們也必須查問應將它們歸於誰。為甚麼連最簡單的題署——如\Quote{達味聖詠}或\Quote{達味作}而無任何附加——也欠缺呢？我的想法是：凡僅有這種題署出現時，所寫的既非聖詠，亦非詩歌，而是某種在聖神引導下的言辭，為能理解的人記錄下來。但有一位希伯來人關於這些末後之事的意見傳到了我這裡；他認為，當有許多聖詠沒有題署，卻以題署為\Quote{達味作}的一篇為首時，所有這些也應算作達味所作的。倘若如此，則那些無題署的聖詠應歸於那些根據題目被正確認為是前面聖詠的作者的人。這本\WorkTitle{聖詠集}也被先知稱為\Quote{聖詠集}，因為據說，在樂器中唯有聖詠琴被敲擊時發出從上而來的聲音，而非像其他樂器那樣從下而來。因此，為使理解它的人熱心實踐此稱呼的類比，並且也向上仰望——從那方向傳來其旋律——他便稱之為\WorkTitle{聖詠集}。因為它完全是至聖聖神的聲音和言辭。\par

3. 讓我們進一步查問，為甚麼有一百五十篇聖詠。五十這個數目是神聖的，這從著名的五旬節慶日便可見出，它象徵脫離勞苦和擁有喜樂。為此，那些日子既不禁食，也不規定屈膝。因為這是為未來時期所保留的大集會的預象。這些時期在以色列地曾有預表，即在希伯來人所稱的\Quote{約貝耳}（喜年）之年，即第五十年，它為奴隸帶來自由，為債務帶來豁免，等等。神聖的福音也知道五十及其同源且相連的數目——即五百——的豁免；因為並非無目的地給了我們五十文和五百文的豁免。因此，為紀念仇敵敗亡並感恩天主的仁慈，頌歌也應包含不只一個五十，而是三個五十，歸於聖父、聖子和聖神之名。\par

4. 再者，五十這個數目包含七個七，或一個安息日的安息日；此外，在這些完整安息日之上，又有一個新的開始，即在八上，一個真正的新安息，超越安息日。凡有能力的人，當在聖詠中以超乎人所能的準確性觀察這點，以便找出每種情況的原因，如同我們將陳述的那樣。例如，第八篇聖詠題署為\Quote{在酒醡上}並非無目的，因為它在八上包含了完美果實；因為享受真葡萄樹果實的時刻不能在八之前。再者，第二篇題署為\Quote{在酒醡上}的是第八十篇，它包含另一個八的數目，即在十倍中。第八十三篇則由兩個神聖數目結合而成：十倍中的八和個位中的三。第五十篇是祈求罪過的赦免和懺悔的祈禱。因為正如按福音，五十獲得赦免，從而確認了對喜年的理解，所以凡以完全懺悔獻上此類懇求的人，希望在除了五十以外的任何數目上獲得赦免。此外，還有一些被稱為\Quote{上升之歌}的，共十五篇，如同聖殿的臺階數，這或許表明\Quote{臺階}包含在七和八這兩個數目內。這些上升之歌始於第一百二十篇之後，該篇僅稱為\Quote{一首聖詠}，正如更準確的抄本所示。這是人生命的完美數目。第一百篇以\BibleQuote{我要歌頌仁慈和審判，上主}開頭，涵蓋了與天主共融的聖人生活。第一百五十篇以\BibleQuote{凡有氣息的，都要讚美上主}結束。\par

5. 但既然，如同我們已經說過，對每一篇這樣做並找出原因非常困難，且非人力所能及，我們便以這些作為綱要自足。只補充一點：涉及歷史材料的聖詠，並非按歷史次序排列。唯一原因在於安排聖詠的數目。例如，第五十一篇的歷史先於第五十篇的歷史。因為眾所周知，厄東人多厄格向撒烏耳誣告達味的事件，先於他與烏黎雅妻子所犯的罪；但並非沒有正當理由而將應在第二的歷史放在第一，因為，如我們先前所說，赦免的主題與五十這個數目有聯繫。因此，凡不配得到赦免者，便越過了五十這個數目，如同厄東人多厄格。因為第五十一篇是論及他的。再者，第三篇也處於同樣情況；它寫於達味逃避其子阿貝沙隆的時候；因此，凡讀\BibleBook{}{列王紀}的人都知道，它應合理地排在第五十一篇和第五十篇之後。\par

凡願意進一步關注這些及類似事項的人，必會自己找到對歷史、題署及聖詠次序的更精確解釋。\par

6. 很可能，关于以下事实也可给出类似的解释：众先知中唯有达味用一件乐器来预言，希腊人称这件乐器为\\OriginalTerm{索尔特里琴}{psaltery}，希伯来人称之为\\OriginalTerm{纳布拉琴}{nabla}，这是唯一完全笔直、没有弯曲的乐器。声音并非来自下部，如同琉特琴和其他某些乐器那样，而是来自上部。因为在琉特琴和里拉琴上，铜片受敲击时从下方发声。但这索尔特里琴其乐音的数字源自上方，为的是我们也能操练寻求上面的事，而不被旋律的快乐拖入肉体的情欲。我认为，这真理也以深远而清晰的方式在乐器构造中向我们预言式地指明，即：那些拥有良好秩序和受训练灵魂的人，有通向上面之事的道路预备好了。再者，一件在上部有悦耳音源的乐器，可被比作基督的身体和祂的圣人们——这唯一保持笔直的乐器；\\BibleQuote{因为祂没有犯过罪，口中也找不到诡诈。}这确实是一件和谐、悦耳、有序的乐器，没有接纳人的不和谐，行事毫不逾矩，而在一切事上维持着向父的和谐；因为祂说：\\BibleQuote{那来自地上的，属于地，也谈论地上的事；那来自天上的，为所见所闻作证。}\par

7. 既然有\\Quote{圣咏}、\\Quote{歌}、\\Quote{歌之圣咏}和\\Quote{圣咏之歌}，我们仍需讨论它们之间的区别。我们认为，\\Quote{圣咏}是指那些单纯用乐器演奏、无人声伴随的（曲调），（它们）是为乐器的音乐旋律而作；那些被称为\\Quote{歌}的，是由人声与音乐协同演唱的；当人声先起，同时乐器也发出适当的音调，和谐地伴奏时，称为\\Quote{歌之圣咏}；而当乐器先起，人声次之，伴随弦乐时，称为\\Quote{圣咏之歌}。就这些术语的字面意义而言，这就是如此。但就神秘释义而言，\\Quote{圣咏}就是通过以善行敲击乐器，即身体，而我们在行动上成功，尽管在默观上并非完全精通；\\Quote{歌}则是当我们脱离实践，默想真理的奥秘并完全认同它们时，对天主和祂的圣言有最高贵的思虑，同时知识光照我们，智慧在我们灵魂中灿烂发光；\\Quote{圣咏之歌}是当善行先起，按照这话：\\BibleQuote{你若渴望智慧，就当遵守诫命，主必将她赐给你。}我们同时领悟智慧，并蒙天主认为配得上明白迄今对我们隐藏的事物真理；\\Quote{歌之圣咏}是当我们以智慧之光默想一些关涉道德的更隐晦问题时，我们首先在行动上变得明智，然后也能说出什么、何时以及如何行动。或许这就是为何最初的标题中从未出现\\Quote{歌}字，而只有\\Quote{圣咏}或\\Quote{圣咏群}；因为圣人不是从默观开始，而是当他按正道教义简单成为信徒后，他就投身于当行之事。也是为此，末后有许多\\Quote{歌}；而凡有\\Quote{登阶}字样的地方，我们找不到\\Quote{圣咏}一词，无论是单独还是附加任何字样，而只有\\Quote{歌}。因为在\\Quote{登阶}（或\\Quote{上行}）中，圣人们将只从事默观。我们按照七十贤士译本释义中的提示所提出的说明，就此主题而言已足够。\par

8. 然而，我们在七十贤士译本、特奥多提翁译本和辛马库译本中，在一些圣咏中——这些还不在少数——发现\\OriginalTerm{希腊词}{\GreekText{δια}}这个词\par

\SourceRecord{Translated by S.D.F. Salmond.；Ante-Nicene Fathers；Vol. 5.；Edited by Alexander Roberts, James Donaldson, and A. Cleveland Coxe.；Buffalo, NY: Christian Literature Publishing Co.,；1886.；<http://www.newadvent.org/fathers/0502.htm>.}
